\documentclass[11pt]{ctexart}
\usepackage[a4paper,margin=1in]{geometry}
\usepackage{graphicx}
\usepackage{xcolor}
\usepackage{hyperref}
\definecolor{refgray}{gray}{0.45}
\title{\textbf{Market Views｜全球投行叙事汇编}\\[0.4em]{\large AI硬件资本开支周期加速与消费电子终端回调形成鲜明分化，中国政策重心转向财政落地与结构性改革，印度即时零售从增量扩张进入存量博弈}}
\author{KC桌面 自动生成}
\date{260718}
\begin{document}
\maketitle
\tableofcontents
\newpage
\section{一页摘要}
\begin{itemize}
  \item 二季度GDP增速放缓至4.3\%，生产与需求分化加剧，全年增速预测下调至4.6\%；政策重心转向加快财政支出与结构性改革，政府债券发行提速是社融主要支撑。
  \item 韩国央行7月全票加息25bp至2.75\%，终端利率路径存在分歧：NOM预测渐进加息至3.25\%，GS预测节奏更快；人民币中间价逆周期因子缓冲收窄至75bp，政策容忍度提升。
  \item 6月社零同比转正至1.0\%，但改善由低基数与一次性因素驱动，消费动能未实质增强；珠宝零售出现K型分化，克重黄金受益于金价下跌，固定定价产品溢价扩大承压。
  \item 中国房地产市场呈现K型分化：一线城市房价企稳但全国销售加速下滑，开发商拿地意愿降至低点；香港地产板块盈利修复启动，但写字楼租金复苏非普涨。
  \item AI基础设施资本开支周期加速，台积电上调2026年营收与资本开支指引，AI需求可见度延伸至2030年；但消费电子终端（PC、智能手机）2026年将面临广泛回调，形成‘AI强、消费弱’格局。
  \item AI数据中心建设从电力获取转向项目落地阶段，管线推进率预期下调；冷却技术路线因单相方案突破面临重新评估，IT预算资金流向高度集中于安全与AI软件。
  \item 中国创新药在ADC和联合疗法领域取得里程碑进展，sac-TMT头对头击败Keytruda+化疗；医疗设备集采政策基调转向合理定价，降价幅度预计比耗材温和。
  \item 印度即时零售行业年化GMV约140亿美元，高速增长阶段已结束，一线城市暗店密度超盈利性均衡点约20\%，行业从增量竞争转向存量竞争，整合与盈利改善成为关键。
\end{itemize}
\section{全球投行叙事汇编}
今日纳入 64 篇投行/券商研究，按 12 个市场、资产与产业主题系统整理。
\newpage
\subsection{宏观与利率：增长动能切换，政策加速落地}
\textbf{二季度GDP增速放缓至4.3\%，生产与需求分化加剧，内需疲软但出口与高技术产业提供支撑。政策重心转向加快财政支出与结构性改革，服务消费与公共融资成为下半年关键变量。}
\subsubsection*{主流共识}
\begin{itemize}
  \item 二季度GDP增速低于预期，全年增速预测下调至4.6\%。
  \item 生产端与需求端分化加剧，工业产出受出口和AI相关产业驱动，内需（零售、投资）持续疲弱。
  \item 下半年财政政策将加速预算落地，而非追加预算，政府债券发行提速是社融主要支撑。
\end{itemize}
\subsubsection*{分歧与条件差异}
\begin{itemize}
  \item 通胀前景：部分机构关注输入性涨价（原油、芯片）推升名义GDP，但内需偏弱下价格效应可持续性存疑。
  \item 信贷修复路径：Citi认为政府债券增量（1.3万亿）难以完全对冲内生信贷缺口（2.2万亿），信用脉冲仍为负；MS则强调公共融资加速的边际改善。
  \item 地缘风险影响：GS区分短期不确定性事件（影响2个季度）与长期结构性结盟（影响14年），后者对GDP的潜在影响更大。
\end{itemize}
\subsubsection*{机构观点}
\paragraph{BARC} 二季度GDP增速放缓至4.3\%，生产与需求分化加剧。需求端指标（零售、投资、出口加权）同比降至-3.6\%，而GDP仍为正增长，背离加深。零售销售6月仅增1\%，固定资产投资连续两月两位数下滑，房地产投资调整24.4\%。工业产出逆势加速至5.3\%，高技术制造业增长14\%，出口韧性是主要支撑。下半年需关注8000亿新融资工具落地及政策会议信号。
{\small\color{refgray}数据：二季度GDP同比4.3\%，低于一季度的5\%\par}
{\small\color{refgray}数据：需求端指标同比-3.6\%，一季度为+3.4\%\par}
{\small\color{refgray}数据：零售销售6月同比+1\%，一季度均值+2.4\%\par}
{\small\color{refgray}数据：固定资产投资6月同比-10\%，房地产投资调整24.4\%\par}
{\small\color{refgray}边际变化：生产与需求背离从季度波动固化为结构性特征，需求端进入负值区间，而工业产出逆势加速。\par}
\paragraph{Citi} 6月信贷数据低于预期，新增人民币贷款1.61万亿（预期1.95万亿），社融3.36万亿（预期3.7万亿）。结构上企业中长期贷款同比少增4580亿，居民贷款均低于去年同期。政府债券融资7690亿为年内低点。下半年政府债券剩余额度7.45万亿，可带来约1.3万亿增量，但内生信贷需求上半年已较去年减少2.2万亿，信用脉冲仍为负。预计7月LPR下调10bp，财政端是社融最确定驱动力。
{\small\color{refgray}数据：新增人民币贷款1.61万亿，低于预期1.95万亿\par}
{\small\color{refgray}数据：社融增量3.36万亿，低于预期3.70万亿\par}
{\small\color{refgray}数据：企业中长期贷款同比少增约4580亿\par}
{\small\color{refgray}数据：政府债券融资7690亿，为年内最低\par}
{\small\color{refgray}边际变化：信贷数据全面低于预期，政府债券增量难以完全对冲内生信贷缺口，信用脉冲下半年仍为负。\par}
\paragraph{MS} 二季度GDP增速4.3\%低于预期，全年预测下调至4.6\%。六月工业增加值反弹至5.3\%，但零售仅增0.2\%，固定资产投资同比-10\%。K型分化加深：高技术制造业受出口支撑，内需走弱。下半年GDP有望回升至4.6\%，前提是预算落地加快及油价正常化。政策重心在加速财政支出而非追加预算，下半年约2万亿存量财政资金可释放。信贷方面，6月广义信贷增速降至7.5\%，但政府债券发行月中已加速至9020亿，超过6月全月。
{\small\color{refgray}数据：全年GDP预测下调至4.6\%（此前4.7\%）\par}
{\small\color{refgray}数据：二季度GDP同比4.3\%，低于市场预期4.5\%\par}
{\small\color{refgray}数据：六月工业增加值同比+5.3\%，零售+0.2\%\par}
{\small\color{refgray}数据：固定资产投资六月同比-10\%\par}
{\small\color{refgray}边际变化：GDP预测下调，但下半年环比增速上调，财政加速落地预期增强；政府债券发行节奏已明显加快。\par}
\paragraph{NOM} 二季度名义GDP增速改善主要由价格效应驱动。制造业名义增速升至8.8\%，但实际增速降至4.8\%，裂创近年之最。出口名义增速20.2\%，同样主要来自价格效应而非量增。金融服务业实际增速6.9\%，但佣金率下降、信贷放缓、净息差收窄至1.40\%，高增长可持续性存疑。房地产与建筑业实际增速-2.3\%，新开工面积-25.7\%。IT服务实际增长10.8\%，AI相关是亮点。
{\small\color{refgray}数据：制造业名义GDP增速8.8\%，实际增速4.8\%\par}
{\small\color{refgray}数据：出口名义增速20.2\%，高于一季度的14.7\%\par}
{\small\color{refgray}数据：金融服务业实际增速6.9\%，名义增速8.9\%\par}
{\small\color{refgray}数据：银行净息差降至1.40\%\par}
{\small\color{refgray}边际变化：名义GDP改善完全由价格效应驱动，实际产出扩张有限；出口强劲主要来自涨价而非量增。\par}
\paragraph{GS} 地缘政治不确定性事件对GDP冲击有限且短暂（约两个季度），但结构性结盟变化影响持续14年。一个标准差的不确定性指数上升使当季GDP下降0.2个百分点，下一季度再降0.1个百分点。2025年上半年全球不确定性指数比2024年下半年高30\%，拖累全球GDP约0.3个百分点。若全球地缘结盟水平回到2010年代中期，到2030年可贡献全球GDP约1\%。新兴市场受影响幅度大于发达市场。
{\small\color{refgray}数据：不确定性指数一个标准差上升，当季GDP下降0.2个百分点\par}
{\small\color{refgray}数据：2025年上半年全球不确定性指数比2024年下半年高30\%\par}
{\small\color{refgray}数据：拖累全球GDP约0.3个百分点\par}
{\small\color{refgray}数据：结构性结盟变化影响持续14年\par}
{\small\color{refgray}边际变化：市场关注短期地缘风险，但长期结构性结盟变化对GDP影响更大且更持久。\par}
\paragraph{JPM} 二季度GDP增速降至4.3\%，全年预测下调至4.6\%。六月工业增加值反弹至5.3\%，但增长集中于出口、设备制造和高技术产业（装备制造+9.3\%，高技术制造+13.3\%）。固定资产投资同比-10\%，房地产投资累计-18\%，私人投资-8.5\%。零售转正至1\%，但汽车销售连续四个月两位数下滑。下半年环比增速上调，财政加速落地是主要支撑。
{\small\color{refgray}数据：全年GDP预测下调至4.6\%（此前4.7\%）\par}
{\small\color{refgray}数据：二季度GDP同比4.3\%\par}
{\small\color{refgray}数据：六月工业增加值同比+5.3\%\par}
{\small\color{refgray}数据：装备制造业+9.3\%，高技术制造业+13.3\%\par}
{\small\color{refgray}边际变化：增长企稳但支撑面收窄，集中于出口和产业升级；下半年环比增速上调，财政加速预期增强。\par}
\paragraph{HSBC} 十五五消费规划明确服务消费为优先方向，目标到2030年服务消费规模达100万亿（年均复合增速4.6\%），高于零售总额60万亿（年均复合增速3.7\%）。重点领域为养老、托育、文旅、健康、体育。短期消费仍疲弱，1-5月零售仅增1.4\%，以旧换新效果递减。政策重心在结构改革而非总量刺激。绿色转型节奏加快，非化石能源占比目标年均提升约1个百分点。
{\small\color{refgray}数据：2030年零售总额目标60万亿，年均复合增速3.7\%\par}
{\small\color{refgray}数据：2030年服务消费规模目标100万亿，年均复合增速4.6\%\par}
{\small\color{refgray}数据：1-5月零售同比+1.4\%\par}
{\small\color{refgray}数据：非化石能源占比目标年均提升约1个百分点\par}
{\small\color{refgray}边际变化：政策重心从商品消费转向服务消费，并设定明确量化目标；绿色转型节奏加快。\par}
\paragraph{GS} 日本通胀叙事从商品驱动转向服务滞后。上调2027财年核心通胀预测至2.3\%，主因日元贬值、存储芯片涨价及食品价格反弹。日元每贬值10日元，新核心CPI约12个月后上升0.25\%，但边际影响递减。日本央行对通胀超标的警惕提高，部分潜在通胀指标已接近2\%。存储芯片价格同比涨7倍，但PC权重仅0.3\%，影响有限；iPhone价格调整（权重1\%）可能贡献0.1\%的CPI影响。
{\small\color{refgray}数据：2027财年核心通胀预测上调至2.3\%\par}
{\small\color{refgray}数据：日元每贬值10日元，新核心CPI上升0.25\%\par}
{\small\color{refgray}数据：存储芯片现货价格同比涨7倍\par}
{\small\color{refgray}数据：PC在CPI中权重0.3\%，影响约0.05\%\par}
{\small\color{refgray}边际变化：通胀结构从商品驱动转向服务滞后，日元贬值边际影响递减但央行关注度上升。\par}
\subsubsection*{关键数据}
\begin{itemize}
  \item 二季度GDP同比4.3\%，低于一季度的5\%
  \item 需求端指标同比-3.6\%，生产端工业产出+5.3\%
  \item 零售销售6月同比+1\%，固定资产投资同比-10\%
  \item 新增人民币贷款1.61万亿，低于预期1.95万亿
  \item 政府债券剩余额度7.45万亿，可带来约1.3万亿增量
  \item 上半年内生信贷需求较去年减少2.2万亿
  \item 制造业名义GDP增速8.8\%，实际增速4.8\%
  \item 2030年服务消费规模目标100万亿，年均复合增速4.6\%
\end{itemize}
\begin{figure}[htbp]
\centering
\includegraphics[width=0.88\linewidth]{figures/fig_001.jpg}
\caption{Figure 1 - \textbackslash{}- Retail sales rose 1\% y/y in June, reversing the 0.6\% decline in May. However, the pace of growth remained well below the average 2.4\% increase recorded in Q1 (Figure 1), reflecting continued headwinds from a challenging labour market, weak private-sector cr}
\end{figure}
\begin{figure}[htbp]
\centering
\includegraphics[width=0.88\linewidth]{figures/fig_026.jpg}
\caption{Figure 1 - Credit and money growth decelerated with the soft numbers. Growth of outstanding TSF decelerated to 7.4\%YoY, a new all-time low. Both M1 and M2 growth softened: M1 growth dipped to 4.0\%YoY as a higher base kicked in (Citi/Mkt: 3.5/4.9\%YoY), a one-year low, and}
\end{figure}
\begin{figure}[htbp]
\centering
\includegraphics[width=0.88\linewidth]{figures/fig_040.jpg}
\caption{Exhibit 1 - Exhibit 1: Marked-to-market cut on full-year GDP, to 4.6\%}
\end{figure}
\begin{figure}[htbp]
\centering
\includegraphics[width=0.88\linewidth]{figures/fig_065.jpg}
\caption{Exhibit 1 - Exhibit 2: Geopolitical Risk and Geopolitical Alignment Are Distinct Phenomena}
\end{figure}
{\small\color{refgray}本节综合 9 篇研究；涉及 BARC、Citi、MS、NOM、GS、HSBC、JPM。}
\newpage
\subsection{外汇与央行：韩元加息路径分化与人民币中间价信号}
\textbf{韩国央行加息周期进入数据依赖阶段，但投行对终端利率和节奏存在分歧；人民币中间价模型显示逆周期因子缓冲收窄，政策容忍度提升；八月季节性指向日元全面走强，高贝塔货币承压。}
\subsubsection*{主流共识}
\begin{itemize}
  \item 韩国央行7月全票加息25bp至2.75\%，符合市场预期，增长因半导体出口显著超预期。
  \item 韩国央行强调未来每次会议均为'现场会议'，政策路径完全取决于数据。
  \item 八月外汇季节性显示日元全面走强，G10货币对日元均贬值，高贝塔和利差货币承压最重。
\end{itemize}
\subsubsection*{分歧与条件差异}
\begin{itemize}
  \item NOM预测韩国央行终端利率3.25\%，维持季度一次渐进加息；GS预测还有两次25bp加息，终端利率3.25\%但节奏可能更快。
  \item NOM认为8月加息概率仅25\%，而市场定价未来一年3-4次加息；GS未排除8月加息可能。
  \item 人民币中间价模型预测6.7739，逆周期因子调整后6.7814，差距收窄至75bp，显示政策对市场定价接受度提高。
\end{itemize}
\subsubsection*{机构观点}
\paragraph{NOM} 韩国央行全票加息25bp至2.75\%，但行长李昌镛拒绝为8月会议预先承诺，强调依赖7月23日Q2 GDP和8月4日7月CPI数据。NOM判断8月加息概率仅25\%，预计维持季度一次渐进加息节奏，终端利率3.25\%。值得关注的是，央行担忧半导体出口带来的收入效应（Q1实际GDI增长13.2\% vs GDP增长3.8\%）可能溢出至国内需求和工资通胀，这是需求侧通胀的新风险。
{\small\color{refgray}数据：韩国央行加息25bp至2.75\%，全票通过\par}
{\small\color{refgray}数据：Q1实际GDP同比+3.8\%，实际GDI同比+13.2\%\par}
{\small\color{refgray}数据：NOM预测终端利率3.25\%，8月加息概率25\%\par}
{\small\color{refgray}数据：市场定价未来一年3-4次加息\par}
{\small\color{refgray}边际变化：韩国央行关注点从供给冲击转向半导体繁荣带来的需求侧通胀风险，收入效应成为新变量。\par}
\paragraph{NOM} 美元/人民币中间价模型预测6.7739，较前次下移170bp，较前一日官方收盘价高61bp。加入逆周期因子后预测6.7814，仅下移95bp，两者差距从约150bp收窄至75bp。差距收窄表明政策端认为市场定价与合意水平偏离缩小，对当前汇率水平的容忍度提高。报告列出7-10月四个关键事件窗口（WAIC大会、政治局会议、主席访美、国庆长假）作为汇率波动催化剂。
{\small\color{refgray}数据：模型预测6.7739，较前次下移170bp\par}
{\small\color{refgray}数据：逆周期因子调整后6.7814，下移95bp\par}
{\small\color{refgray}数据：差距从150bp收窄至75bp\par}
{\small\color{refgray}数据：关键事件窗口：7月WAIC、7月底政治局会议、9月24日主席访美、10月初国庆长假\par}
{\small\color{refgray}边际变化：逆周期因子缓冲幅度收窄，政策对市场定价的接受度提高，中间价定价机制透明度上升。\par}
\paragraph{DB} 八月外汇季节性主线是日元全面走强，而非美元单向升值。基于60个货币对、二十余年日收盘价和交易时段数据，G10货币对日元均贬值，统计显著性高于对美元贬值：英镑/日元t统计量-2.9，澳元/日元-2.6，新西兰元/日元-2.5。价格发现集中在伦敦-纽约重叠时段，表明机构资金和避险行为是主要驱动力。高贝塔货币（ZAR、MXN、AUD、NZD）承压最重，新西兰元还受南半球冬季乳制品出口淡季拖累。
{\small\color{refgray}数据：G10货币对日元8月均贬值，英镑/日元t统计量-2.9\par}
{\small\color{refgray}数据：南非兰特/日元t统计量-3.0，墨西哥比索/日元-2.7\par}
{\small\color{refgray}数据：价格发现集中在伦敦-纽约重叠时段\par}
{\small\color{refgray}数据：新西兰元受南半球冬季乳制品出口淡季额外拖累\par}
{\small\color{refgray}边际变化：八月季节性驱动机制从年末流动性对冲转向信息敏感型避险行为，日元升值的价格发现集中在信息密度最高的交易窗口。\par}
\paragraph{GS} 韩国央行加息25bp至2.75\%，声明指出2026年实际GDP增长将'显著超过'5月预测的2.6\%，半导体出口是核心引擎。通胀预计与2.7\%预测一致，但核心通胀可能高于2.4\%预测，不确定性来自成本端和汇率。家庭贷款和房价加速上涨，金融稳定风险上升。行长明确每次会议均为'现场会议'，GS预计还有两次25bp加息（Q4和Q1），终端利率3.25\%。
{\small\color{refgray}数据：韩国央行加息25bp至2.75\%，全票通过\par}
{\small\color{refgray}数据：2026年GDP增长将'显著超过'5月预测的2.6\%\par}
{\small\color{refgray}数据：2026年整体通胀预测2.7\%，核心通胀可能高于2.4\%\par}
{\small\color{refgray}数据：GS预计还有两次25bp加息，终端利率3.25\%\par}
{\small\color{refgray}边际变化：韩国央行从明确加息路径转向完全数据依赖，增长预期上调幅度超市场预期，但半导体持续性不确定性增加。\par}
\subsubsection*{关键数据}
\begin{itemize}
  \item 韩国央行加息25bp至2.75\%，全票通过
  \item NOM预测终端利率3.25\%，8月加息概率25\%
  \item GS预计还有两次25bp加息，终端利率3.25\%
  \item 人民币中间价模型预测6.7739，逆周期因子调整后6.7814，差距收窄至75bp
  \item Q1韩国实际GDP同比+3.8\%，实际GDI同比+13.2\%
  \item G10货币对日元8月均贬值，英镑/日元t统计量-2.9
  \item 南非兰特/日元t统计量-3.0
  \item 韩国2026年GDP增长将'显著超过'5月预测的2.6\%
\end{itemize}
\begin{figure}[htbp]
\centering
\includegraphics[width=0.88\linewidth]{figures/fig_053.jpg}
\caption{Figure 1 - Figure 1: USD/JPY exhibits strong seasonality in August}
\end{figure}
\begin{figure}[htbp]
\centering
\includegraphics[width=0.88\linewidth]{figures/fig_054.jpg}
\caption{Figure 2 - Figure 2: Breadth of seasonality highest in August (based on London closing prices)}
\end{figure}
\begin{figure}[htbp]
\centering
\includegraphics[width=0.88\linewidth]{figures/fig_056.jpg}
\caption{Figure 4 - Figure 4 - Figure 8). Figure 4: August a risk-off month - G10 and EM currencies weakening versus USD and JPY}
\end{figure}
\begin{figure}[htbp]
\centering
\includegraphics[width=0.88\linewidth]{figures/fig_055.jpg}
\caption{Figure 2 - Figure 4 - Figure 8).}
\end{figure}
{\small\color{refgray}本节综合 4 篇研究；涉及 NOM、DB、GS。}
\newpage
\subsection{中国消费与零售：6月数据改善，但结构分化与成本压力并存}
\textbf{6月社零同比转正至1.0\%，略超预期，但改善主要由低基数和一次性因素（618、补贴）驱动，两年复合增速持平于5月。消费动能未实质增强，且CPI-PPI剪刀差扩大至-3.1\%，成本压力正在侵蚀下游利润。珠宝零售出现K型分化：克重黄金受益于金价下跌需求释放，固定定价产品溢价扩大承压。运动品牌中，李宁库存改善至4个月，但零售流水仍承压。预计二季度业绩期更多公司将下调全年指引。}
\subsubsection*{主流共识}
\begin{itemize}
  \item 6月社零同比转正至1.0\%，高于预期的-0.1\%，但两年复合增速5.8\%与5月持平，改善主要来自低基数，而非消费动能实质增强。
  \item 消费结构分化明显：化妆品（+12.6\%）、烟酒（+12.1\%）、通讯器材（+16.5\%）表现突出，但汽车（-16.1\%）、家电（-8.7\%）、家具（-6.6\%）持续疲弱。
  \item 珠宝零售中，克重黄金产品受益于6月金价环比下跌11\%，需求集中释放；固定定价产品因溢价扩大至74\%而承压。
\end{itemize}
\subsubsection*{分歧与条件差异}
\begin{itemize}
  \item 关于消费复苏节奏：Citi认为6月数据改善叠加政策获批，可能触发短期情绪修复；JPM则认为二季度零售总额同比仅增0.2\%，较一季度的2.4\%明显放缓，消费动能仍在走弱。
  \item 关于珠宝品牌前景：Citi认为老铺黄金二季度可能是销售和盈利底部，估值有吸引力（8.7倍PE）；GS则强调六福内地关店150家远超预期，渠道优化加速，但固定定价珠宝下滑7\%显示品牌溢价能力减弱。
  \item 关于运动品牌：Citi对李宁开启短期正面观察，认为悲观预期已充分定价；JPM指出李宁二季度零售流水同比下滑低个位数，但库存改善至4个月，管理层维持全年指引。
\end{itemize}
\subsubsection*{机构观点}
\paragraph{Citi} 6月社零同比增速从5月的-0.6\%回升至+1.0\%，餐饮收入升至+1.2\%，方向性变化叠加促消费政策获批，可能触发短期情绪修复。珠宝零售中，6月金价环比下跌11\%触发克重黄金需求集中释放，周大福内地同店增长20\%（克重黄金+35\%），六福内地同店增速超20\%；但老铺黄金固定定价产品溢价扩大至74\%，二季度业绩存在不确定性，但可能是阶段性底部。重申对李宁和海底捞的短期正面观察，李宁盈利预测比共识低8-10\%，悲观预期已充分定价。
{\small\color{refgray}数据：6月社零同比+1.0\%（5月-0.6\%），餐饮+1.2\%\par}
{\small\color{refgray}数据：6月金价环比-11\%，金银珠宝零售同比-3.4\%（4月-21.3\%，5月-8.9\%）\par}
{\small\color{refgray}数据：老铺黄金产品溢价率扩大至74\%（此前约50\%）\par}
{\small\color{refgray}数据：周大福4-5月内地同店+20\%，克重黄金+35\%\par}
{\small\color{refgray}边际变化：此前市场对消费板块悲观预期过度，6月数据改善和政策信号可能触发短期情绪修复；珠宝零售中克重黄金需求在金价下跌后集中释放，而固定定价产品溢价被动扩大承压。\par}
\paragraph{MS} 6月社零同比+1.0\%，高于市场预期的-0.1\%，但剔除汽车后增速3.0\%，以2019年为基准的年复合增速仅从3.2\%微升至3.4\%，修复力度有限。化妆品同比+12.6\%受618和补贴驱动，烟酒+12.1\%受低基数影响，两者均不宜外推为消费意愿系统性回升。必选消费全面改善，可选消费跌幅收窄但未转正。线上零售增速放缓至2.0\%（5月3.4\%），部分因618促销提前启动。
{\small\color{refgray}数据：6月社零同比+1.0\%，高于预期的-0.1\%\par}
{\small\color{refgray}数据：剔除汽车后零售增速+3.0\%\par}
{\small\color{refgray}数据：化妆品同比+12.6\%（5月+2.5\%），烟酒+12.1\%（5月+4.8\%）\par}
{\small\color{refgray}数据：线上零售同比+2.0\%（5月+3.4\%）\par}
{\small\color{refgray}边际变化：6月数据超预期转正，但MS强调修复路径不平坦，结构分化是核心特征；化妆品和烟酒的高增长分别归因于一次性因素和低基数，不宜直接外推。\par}
\paragraph{GS} 六福第一财季港澳同店销售增长41\%，内地自营门店增长16\%，黄金产品贡献绝大部分增量。但内地门店净关闭150家，远超市场此前对半个财年150家的预期，渠道优化加速。固定价格珠宝在内地自营渠道同比下滑7\%，品牌溢价变现能力减弱。港澳市场整体同店销售恢复至2018年同期78\%，黄金产品已完全恢复（103\%），但固定价格珠宝仅恢复至48\%。
{\small\color{refgray}数据：六福港澳同店销售增长41\%\par}
{\small\color{refgray}数据：内地自营门店同店增长16\%，加盟店+25\%\par}
{\small\color{refgray}数据：内地净关闭门店150家（上季度71家），远超预期\par}
{\small\color{refgray}数据：固定价格珠宝内地自营同比-7\%\par}
{\small\color{refgray}边际变化：六福内地关店速度远超预期，单季度净关闭150家（此前预期整个上半财年150家），渠道优化加速；黄金产品增长强劲，但固定定价珠宝下滑显示消费者更倾向保值型消费。\par}
\paragraph{JPM} 6月社零同比转正至1.0\%，但两年复合增速5.8\%与5月持平，改善主要来自低基数。二季度零售总额同比仅增0.2\%，较一季度的2.4\%明显放缓。CPI-PPI剪刀差扩大至-3.1\%，成本压力向下游传导，多数消费企业定价能力不足。预计更多公司将在二季度业绩期下调全年指引。李宁二季度零售流水同比下滑低个位数，但库存从约5个月改善至约4个月，库龄结构健康，管理层维持全年指引，库里合作费用摊销对2026年影响有限。
{\small\color{refgray}数据：6月社零同比+1.0\%，两年复合增速5.8\%\par}
{\small\color{refgray}数据：二季度零售总额同比+0.2\%（一季度+2.4\%）\par}
{\small\color{refgray}数据：CPI-PPI剪刀差-3.1\%\par}
{\small\color{refgray}数据：李宁库存从约5个月改善至约4个月\par}
{\small\color{refgray}边际变化：JPM强调二季度消费动能较一季度明显放缓，6月改善主要来自低基数；李宁库存改善但零售流水承压，管理层维持全年指引，市场对库里合作一次性拉低利润的担忧可能被高估。\par}
\subsubsection*{关键数据}
\begin{itemize}
  \item 6月社零同比+1.0\%，高于预期的-0.1\%，两年复合增速5.8\%
  \item 二季度零售总额同比+0.2\%，较一季度的+2.4\%放缓
  \item CPI-PPI剪刀差扩大至-3.1\%
  \item 6月金价环比-11\%，金银珠宝零售同比-3.4\%（4月-21.3\%，5月-8.9\%）
  \item 老铺黄金产品溢价率扩大至74\%
  \item 六福内地单季度净关闭门店150家，远超预期
  \item 李宁库存从约5个月改善至约4个月
  \item 化妆品同比+12.6\%（618+补贴），烟酒+12.1\%（低基数）
\end{itemize}
\begin{figure}[htbp]
\centering
\includegraphics[width=0.88\linewidth]{figures/fig_031.jpg}
\caption{Figure 1 - Jun 26 note: Luk Fook Holdings (International) (0590.HK) - FY27E Growth Target Stronger Than Expected. Figure 1. Price comparison of Laopu and major traditional gold jeweler © 2026 Citi Inc. No redistribution without Citi's written permission. Figure 2. Laopu'}
\end{figure}
\begin{figure}[htbp]
\centering
\includegraphics[width=0.88\linewidth]{figures/fig_033.jpg}
\caption{Figure 1 - Brian Cho Figure 1. YoY change of China total restaurant sales © 2026 Citi Inc. No redistribution without Citi's written permission. \textbackslash{}*Jan-Feb data were reported on combined basis Figure 2. YoY change of China total consumer goods retail sales © 2026 Citi Inc.}
\end{figure}
\begin{figure}[htbp]
\centering
\includegraphics[width=0.88\linewidth]{figures/fig_096.jpg}
\caption{Exhibit 1 - Exhibit 3: SSS of gold products accelerated in both ML China and HK\&Macau in Jun-quarter.}
\end{figure}
\begin{figure}[htbp]
\centering
\includegraphics[width=0.88\linewidth]{figures/fig_115.jpg}
\caption{Figure 1 - Figure 1: China retail sales YoY trend Figure 2: China retail sales YoY trend – key categories Figure 3: Retail sales YoY growth ranked by category: April 2026}
\end{figure}
{\small\color{refgray}本节综合 6 篇研究；涉及 Citi、MS、GS、JPM。}
\newpage
\subsection{中国地产与基建：分化加剧，盈利修复启动}
\textbf{中国房地产市场呈现K型分化：一线城市房价企稳但全国销售加速下滑，开发商拿地意愿降至低点。香港地产板块则开启盈利上行周期，写字楼租金回暖，但区域分化显著。政策定位从增长引擎转向消费载体，但市场出清仍需时间。}
\subsubsection*{主流共识}
\begin{itemize}
  \item 全国商品房销售面积和金额同比降幅在6月扩大，市场回暖未能持续。
  \item 开发商拿地意愿低迷，上半年百强房企拿地金额同比降约40\%，供给端调整周期长于需求端。
  \item 一线城市房价相对企稳，但全国70城整体仍处于调整通道，二手房价格跌幅扩大。
  \item 香港地产板块盈利修复启动，驱动力来自融资成本下降和利润率改善，而非收入增长。
\end{itemize}
\subsubsection*{分歧与条件差异}
\begin{itemize}
  \item 一线城市二手房价格走势存在数据分歧：NBS口径连续4个月环比正增长，但中原指数6月环比转负，深圳和广州实际成交价可能更弱。
  \item 中海地产与华润置地同为国企龙头，但利润前景分化：中海利润率承压但减值风险被高估，华润置地则依赖REIT处置收益和租金收入缓冲开发业务减值。
  \item 香港写字楼市场复苏并非普涨：中环租金环比涨3.3\%，但港岛东和九龙东仍在下跌，核心与非核心地段租金差距拉大。
  \item 政策定位从“扩张供给”转向“提升品质”，但市场量价调整仍在深化，政策意图与市场现实之间存在张力。
\end{itemize}
\subsubsection*{机构观点}
\paragraph{GS} 中海地产上半年合约销售同比增长12\%至1340亿元，显著跑赢国企同行，但毛利率承压，预计2026年毛利率约14\%。减值风险被市场高估，公司超80\%土储位于一二线城市，累计减值比例在国企中最低。拿地节奏主动放缓，上半年土地投资仅占合同销售的高个位数百分比，新项目毛利率约29\%。预计2026年核心每股收益下降6\%，2027年起恢复增长。
{\small\color{refgray}数据：中海地产1H26合约销售同比+12\%至1340亿元\par}
{\small\color{refgray}数据：GS预计2026年毛利率约14\%，处于管理层指引13\%-15\%中枢\par}
{\small\color{refgray}数据：超80\%土储位于一二线城市，累计减值比例在国企中最低\par}
{\small\color{refgray}数据：上半年拿地强度为高个位数百分比，远低于2025年全年约50\%\par}
{\small\color{refgray}边际变化：GS将2026-2028年毛利率预测下调约1个百分点，但强调减值不确定性被高估，市场已计入同行中最高的减值损失。\par}
\paragraph{GS} 华润置地上半年净利润预计同比持平，开发业务毛利率预测下调4个百分点至13.4\%，低线城市老库存减值规模可能超15.5亿元。但租金收入同比增长13\%，成都万象城REIT上市提供关键缓冲。合同销售额同比增长6\%至1170亿元，拿地强度约26\%，新项目毛利率约24\%。
{\small\color{refgray}数据：华润置地1H26净利润预计同比持平\par}
{\small\color{refgray}数据：开发业务毛利率预测下调4个百分点至13.4\%\par}
{\small\color{refgray}数据：低线城市老库存减值规模可能超15.5亿元\par}
{\small\color{refgray}数据：租金收入同比+13\%，明显高于同业个位数增速\par}
{\small\color{refgray}边际变化：GS将2026年开发业务毛利率预测下调4个百分点，并调降2027-2028年预期，反映低线城市房价调整对利润表的实质性影响。\par}
\paragraph{MS} 6月全国商品房销售额和销售面积同比分别下降13.9\%和14.3\%，降幅较5月扩大。二手房市场明显走弱，高能级城市实时二手房销售增速从5月的26\%降至6月的10\%。开发商拿地意愿降至低点，上半年百强房企拿地面积同比降约40\%。三季度二手房销售可能同比转负，新房销售在低基数下仍将进一步下滑。
{\small\color{refgray}数据：6月全国商品房销售额同比-13.9\%，销售面积同比-14.3\%\par}
{\small\color{refgray}数据：高能级城市实时二手房销售增速从5月的26\%降至6月的10\%\par}
{\small\color{refgray}数据：上半年百强房企拿地面积同比-40\%\par}
{\small\color{refgray}数据：6月新开工面积同比-26\%，竣工面积同比-25\%\par}
{\small\color{refgray}边际变化：五月短暂回暖未能持续，销售降幅扩大；二手房先行指标走弱，三季度可能转负。\par}
\paragraph{MS} 香港写字楼市场回暖，2Q26整体租金环比涨1.7\%，中环核心区环比涨3.3\%，年初至今累计涨7.3\%。空置率改善0.4个百分点至13.1\%，中环空置率降至8.8\%。但复苏并非普涨，港岛东和九龙东租金仍在环比下跌。零售租金同比降8.4\%，但降幅较一季度的-10\%收窄。甲级写字楼资本价值环比微涨0.8\%，但仍较2018年峰值低约50\%。
{\small\color{refgray}数据：2Q26香港整体写字楼租金环比+1.7\%，中环环比+3.3\%\par}
{\small\color{refgray}数据：中环空置率降至8.8\%，环比改善0.8个百分点\par}
{\small\color{refgray}数据：港岛东和九龙东租金环比-0.8\%和-1.1\%\par}
{\small\color{refgray}数据：优质零售租金同比-8.4\%，好于一季度的-10\%\par}
{\small\color{refgray}边际变化：写字楼租金趋势从下跌转为上涨，中环领涨；零售租金跌幅收窄，边际改善信号出现。\par}
\paragraph{JPM} 6月一线城市房价连续第四个月环比正增长，但全国70城整体趋势分化。NBS口径一线二手房环比+0.3\%，但中原指数环比-0.2\%，深圳和广州实际成交价可能更弱。全国住宅销售额同比降幅从5月的8\%扩大至6月的14\%，全年销售预测从-7\%下调至-10\%。供给调整（新开工同比-26\%）为房价企稳创造条件，但需求端能否跟上存疑。
{\small\color{refgray}数据：6月一线城市二手房NBS口径环比+0.3\%，连续4个月正增长\par}
{\small\color{refgray}数据：中原指数6月环比-0.2\%，深圳-0.8\%拖累\par}
{\small\color{refgray}数据：全国住宅销售额6月同比-14\%，较5月的-8\%扩大\par}
{\small\color{refgray}数据：全年销售预测从-7\%下调至-10\%\par}
{\small\color{refgray}边际变化：一线城市二手房价格出现数据分歧，NBS与中原指数背离；全国销售下滑加速，全年预测下调。\par}
\paragraph{JPM} 香港地产板块开启跨年度盈利上行周期，2026年整体盈利增速约9\%，股息增长约2-3\%。驱动力来自融资成本下降和利润率改善，而非收入增长。子板块动量排序：中环写字楼>住宅=内地零售>非核心写字楼>内地写字楼。恒基地产预计增长42\%（高利润项目入账），新鸿基地产预计增长5\%，长实集团剔除一次性收益后核心利润下降6\%。
{\small\color{refgray}数据：2026年香港地产整体盈利增速约9\%，股息增长约2-3\%\par}
{\small\color{refgray}数据：恒基地产预计利润+42\%，受益于高利润项目“The Legacy”\par}
{\small\color{refgray}数据：新鸿基地产预计+5\%，存在上行空间\par}
{\small\color{refgray}数据：长实集团剔除一次性收益后核心利润-6\%\par}
{\small\color{refgray}边际变化：香港地产板块从2-3年盈利调整周期中走出，盈利上行周期开启，但修复主要来自成本端而非收入端。\par}
\paragraph{JPM} 住房政策定位从增长引擎转向消费载体，与汽车、家电并列在居民消费篮子中。政策重心从“让更多人买房”转向“提升居住品质”，包括保障房、改善性住房供给、发展租赁市场。但市场现实与政策意图存在张力：6月住房活动指数继续调整，新房销售面积同比-15.3\%，二手房价格跌幅扩大至-0.32\%。新房价格较2021年峰值累计下跌13.9\%。
{\small\color{refgray}数据：6月住宅新开工面积同比-26.6\%，较5月的-22.9\%进一步扩大\par}
{\small\color{refgray}数据：6月商品房销售面积同比-15.3\%，较5月的-12.4\%扩大\par}
{\small\color{refgray}数据：70城二手房价格环比-0.32\%，较5月的-0.26\%扩大\par}
{\small\color{refgray}数据：新房价格较2021年峰值累计下跌13.9\%\par}
{\small\color{refgray}边际变化：政策框架从“扩张供给”转向“提升品质”，但量价调整仍在深化，此前被视为积极信号的复苏势头减弱。\par}
\subsubsection*{关键数据}
\begin{itemize}
  \item 中海地产1H26合约销售同比+12\%至1340亿元
  \item 华润置地低线城市老库存减值规模可能超15.5亿元
  \item 上半年百强房企拿地面积同比-40\%
  \item 6月全国商品房销售额同比-13.9\%
  \item 一线城市二手房NBS口径环比+0.3\%，中原指数环比-0.2\%
  \item 2Q26香港中环写字楼租金环比+3.3\%，年初至今累计+7.3\%
  \item 2026年香港地产整体盈利增速约9\%
  \item 新房价格较2021年峰值累计下跌13.9\%
\end{itemize}
\begin{figure}[htbp]
\centering
\includegraphics[width=0.88\linewidth]{figures/fig_120.jpg}
\caption{Figure 1 - Figure 1: Tier-1 cities – secondary home price index (NBS and Centraline) with major nationwide easing/narrative change Figure 2: Tier-1 cities – secondary home price M/M growth (NBS vs. Centraline)}
\end{figure}
\begin{figure}[htbp]
\centering
\includegraphics[width=0.88\linewidth]{figures/fig_046.jpg}
\caption{Exhibit 1 - Exhibit 1: Office Rent Index - YoY Change Exhibit 2: Retail Rent Index - YoY Change MS ASIA LIMITED+ Praveen K Choudhary}
\end{figure}
\begin{figure}[htbp]
\centering
\includegraphics[width=0.88\linewidth]{figures/fig_130.jpg}
\caption{Figure 1 - Figure 2: 1H26E / FY26E / 1HFY27E DPS Y/Y growth Core net profit Y/Y Figure 3: Semi-annual core net profit Y/Y growth Note: including companies reporting 1H26 results only}
\end{figure}
\begin{figure}[htbp]
\centering
\includegraphics[width=0.88\linewidth]{figures/fig_122.jpg}
\caption{Figure 3 - Figure 3: 70-city primary and secondary home price M/M trend Note: nationwide easing is marked as green. Figure 4: NBS 70-city home price index – primary M/M growth by tier Figure 5: NBS 70-city price index – \% of cities with}
\end{figure}
{\small\color{refgray}本节综合 7 篇研究；涉及 GS、MS、JPM。}
\newpage
\subsection{AI硬件升级与终端需求分化：科技硬件的结构性机会与风险}
\textbf{AI基础设施的资本开支周期正在加速，台积电将2026年营收增速上调至略超40\%，资本开支指引提升至600-640亿美元，且管理层将AI需求可见度延伸至2030年。但与此同时，消费电子终端（PC、智能手机）在2026年将面临广泛回调，PC出货量预计下滑14\%，智能手机下滑11.3\%，形成‘AI强、消费弱’的分化格局。中国半导体进口数据显示高端芯片需求韧性超预期，6月进口价值同比增72.3\%，均价跳升61.6\%，进一步验证AI与智能驾驶的结构性拉动。}
\subsubsection*{主流共识}
\begin{itemize}
  \item AI基础设施资本开支周期加速，台积电上调2026年营收与资本开支指引，AI需求可见度延伸至2030年。
  \item 2026年消费电子终端（PC、智能手机）出货量将出现广泛回调，PC出货量预计下滑14\%，智能手机下滑11.3\%。
  \item 中国半导体进口高端化趋势明确，进口均价连续两个月同比涨超60\%，AI与智能驾驶是主要驱动力。
\end{itemize}
\subsubsection*{分歧与条件差异}
\begin{itemize}
  \item AI服务器升级路径分歧：MS强调GPU（Vera Rubin/Kyber）与ASIC（TPU/Trainium）双线设计升级带来的供应链价值提升，而GS更聚焦台积电先进制程与CoWoS产能扩张对设备商的拉动。
  \item PC市场量价背离：GS认为AI PC因价格敏感度低、边缘计算需求上升，渗透率将从36\%跃升至59\%，但JPM预测整体PC出货量2026年下滑8.5\%，显示对AI PC能否抵消总量下滑存在分歧。
  \item 设备涨价持续性：GS日本设备团队认为台积电将通胀列为资本开支上调因素，设备涨价正从个案变为趋势；但JPM对元器件库存高企（62天）表示担忧，认为产能利用率面临不确定性。
\end{itemize}
\subsubsection*{机构观点}
\paragraph{MS} 大中华科技硬件供应链将在2026H2-2027年迎来由GPU和ASIC服务器机架设计升级驱动的结构性变化。Vera Rubin平台、Kyber架构及LPX/Vera CPU机架是未来两年最明确的设计变更锚点，800V直流电源方案、PCB/ABF载板、数据中心网络互联/CPO升级是短期催化剂。AI ASIC服务器（TPU、Trainium）的规模扩张独立于GPU路线，要求ODM和零部件供应商同时服务两套技术标准。消费电子需求疲软反而使市场更聚焦AI服务器这一确定性增量方向。
{\small\color{refgray}数据：800V直流电源方案被列为近端报价上行空间可见的领域\par}
{\small\color{refgray}数据：AI ASIC服务器（TPU、Trainium平台）量价齐升\par}
{\small\color{refgray}数据：GB300服务器机柜交付顺利\par}
{\small\color{refgray}边际变化：从关注AI需求总量转向聚焦设计升级带来的单机价值提升和技术迭代，而非简单的出货量扩张。\par}
\paragraph{GS} 全球PC市场2026年将经历量价背离：出货量同比下滑14\%至2.55亿台，但ASP同比上涨10\%至1025美元。AI PC凭借对价格上涨的低敏感性和边缘计算需求，出货量预计增长39\%至1.5亿台，渗透率从36\%跃升至59\%。成本端，内存和CPU价格上升是出货量下修的直接原因，但AI PC用户对价格上涨的容忍度更高。品牌格局变化中，Apple份额从9\%升至14\%，Lenovo份额持续提升至31\%。
{\small\color{refgray}数据：2026年全球PC出货量2.55亿台，同比-14\%\par}
{\small\color{refgray}数据：2026年AI PC出货量1.5亿台，渗透率59\%\par}
{\small\color{refgray}数据：2026年PC ASP 1025美元，同比+10\%\par}
{\small\color{refgray}数据：Apple PC份额从2025年9\%升至2028年14\%\par}
{\small\color{refgray}边际变化：此前GS预测2026年PC出货量下滑10\%、2027年恢复增长3\%，现大幅下修至2026年-14\%、2027年-5\%。\par}
\paragraph{GS} 中国半导体进口数据显示高端芯片需求韧性超预期。6月IC进口价值同比增72.3\%，进口量仅增6.6\%，进口均价同比跳升61.6\%，连续两个月涨超60\%。5月半导体行业总营收达320亿美元，同比增87.2\%，IC产量同比增22.9\%，连续三个月加速。电子制造业库存天数62天，高于过去几年同期水平。生成式AI和ADAS/自动驾驶是背后的结构性推力。设备采购结构分化，测试设备进口额同比增50\%，光刻机进口量降40\%但均价涨64\%。
{\small\color{refgray}数据：6月中国IC进口价值同比+72.3\%\par}
{\small\color{refgray}数据：6月IC进口均价同比+61.6\%\par}
{\small\color{refgray}数据：5月半导体行业总营收320亿美元，同比+87.2\%\par}
{\small\color{refgray}数据：5月IC产量同比+22.9\%\par}
{\small\color{refgray}边际变化：进口均价从2026年初的同比+28.2\%加速至5-6月的+69.7\%/+61.6\%，显示高端化趋势显著强化。\par}
\paragraph{GS} 台积电7月法说会将2026年营收增速指引从‘高于30\%’上调至‘略高于40\%’，资本开支指引从520-560亿美元提升至600-640亿美元。管理层首次明确将通胀导致的采购成本上升列为资本开支上调因素，对日本半导体设备行业具有结构性含义——设备涨价正从个案讨论变成行业趋势。台积电同时确认A14节点按计划2028年量产，先进制程设备需求周期至少还有两到三年可见度。日本设备商（Lasertec、Ebara、Disco、东京电子）因对台积电暴露度高而直接受益。
{\small\color{refgray}数据：台积电2026年营收增速指引从>30\%上调至略超40\%\par}
{\small\color{refgray}数据：资本开支指引从520-560亿美元提升至600-640亿美元\par}
{\small\color{refgray}数据：台积电在亚利桑那追加1000亿美元投资\par}
{\small\color{refgray}数据：A14节点按计划2028年量产\par}
{\small\color{refgray}边际变化：台积电首次将通胀列为资本开支上调的独立变量，设备涨价逻辑从材料成本传导升级为客户接受设备定价上调。\par}
\paragraph{JPM} 全球电子元器件行业2025H2-2026年多数终端市场将出现同比增速转负。智能手机2026年出货量预计同比下滑11.34\%，降幅超过2022年；PC出货量2026年下滑8.53\%。服务器是相对亮点，2026年出货量预计增长14.5\%，其中AI服务器2025年增长近70\%，2026年增速放缓至21.3\%。汽车市场中xEV持续渗透，HEV增长稳定。日本14家核心元器件制造商的库存指标显示，终端需求的真实强度与元器件厂商的产能规划之间存在扩大的时间差。
{\small\color{refgray}数据：2026年全球智能手机出货量同比-11.34\%\par}
{\small\color{refgray}数据：2026年全球PC出货量同比-8.53\%\par}
{\small\color{refgray}数据：2026年全球服务器出货量同比+14.5\%\par}
{\small\color{refgray}数据：AI服务器2025年出货量同比+70\%，2026年增速放缓至+21.3\%\par}
{\small\color{refgray}边际变化：此前市场普遍期待的‘AI拉动全面复苏’需要重新评估，消费电子与AI服务器需求分化加剧。\par}
\subsubsection*{关键数据}
\begin{itemize}
  \item 台积电2026年营收增速指引从>30\%上调至略超40\%，资本开支指引提升至600-640亿美元
  \item 2026年全球PC出货量同比-14\%至2.55亿台，AI PC渗透率从36\%升至59\%
  \item 6月中国IC进口价值同比+72.3\%，进口均价同比+61.6\%
  \item 5月中国半导体行业总营收320亿美元，同比+87.2\%
  \item 2026年全球智能手机出货量预计同比-11.34\%
  \item AI服务器2025年出货量同比+70\%，2026年增速放缓至+21.3\%
  \item 台积电A14节点按计划2028年量产
  \item 电子制造业库存天数62天，高于过去几年同期水平
\end{itemize}
\begin{figure}[htbp]
\centering
\includegraphics[width=0.88\linewidth]{figures/fig_099.jpg}
\caption{Exhibit 1 - Exhibit 1: TSMC financial snapshot Exhibit 2: TSMC's revenue growth outlook Exhibit 3: TSMC's margin outlook Exhibit 4: TSMC's annual CoWoS shipment growth trend}
\end{figure}
\begin{figure}[htbp]
\centering
\includegraphics[width=0.88\linewidth]{figures/fig_076.jpg}
\caption{Exhibit 2 - Exhibit 2: Global PC Shipments: -14\%/-5\%/0\% YoY in 2026-28E Exhibit 3: Global PC Revenues: -5\%/ -2\%/ +3\% YoY in 2026 / 27E Exhibit 4: PC ASP: increasing pricing due to specification upgrades and rising BoM}
\end{figure}
\begin{figure}[htbp]
\centering
\includegraphics[width=0.88\linewidth]{figures/fig_082.jpg}
\caption{Exhibit 3 - Exhibit 3: IC import value was +72.3\% YoY in Jun 2026 (vs. +68.0\% YoY in May 2026) Exhibit 4: IC imports volume at -1.0\% / +6.6\% YoY in May / Jun 2026 \#\# Exhibit 5: May 2026 inventory days were at 62 vs. 58/56/60 days in May 20}
\end{figure}
\begin{figure}[htbp]
\centering
\includegraphics[width=0.88\linewidth]{figures/fig_090.jpg}
\caption{Exhibit 1 - Exhibit 1: Sales exposure to TSMC (based on FY25 results) Lasertec and Ulvac figures are based on our estimates. Exhibit 2: TP risks and methodology for companies mentioned \textbackslash{}*=on Conviction List}
\end{figure}
{\small\color{refgray}本节综合 6 篇研究；涉及 MS、GS、JPM。}
\newpage
\subsection{AI基础设施与数据中心：建设落地分化，冷却技术路线存变数}
\textbf{AI数据中心建设正从电力获取转向项目落地阶段，管线推进率预期下调；冷却技术路线因单相方案突破面临重新评估；IT预算温和改善但资金流向高度集中，安全与AI软件为明确赢家。}
\subsubsection*{主流共识}
\begin{itemize}
  \item 数据中心项目落地难度上升，公众反对（电价、用水、就业）成为主要障碍，开发商采取多地点并行申请策略对冲不确定性。
  \item 电力设备交期仍紧张（燃气轮机排至2030年），但涨价幅度从30-50\%降至温和通胀，产能扩张与二手设备释放缓解供给压力。
  \item 2026年全球IT预算增速预期小幅回升至+3.8\%，软件仍是增长最快板块（+4.1\%），安全软件以+12.4\%增速领跑。
\end{itemize}
\subsubsection*{分歧与条件差异}
\begin{itemize}
  \item 数据中心冷却技术路线：单相DTC能否覆盖下一代GPU散热需求存在分歧。CoolIT 15kW冷板证明单相方案仍有显著余量，但行业主流仍预期两相DTC将在12-18个月内商业化。
  \item AI支出资金来源：增量资金占比下降（从61\%降至54\%），重新分配占比上升（从23\%升至32\%），部分CIO开始从现有软件预算中调配资源支持AI。
  \item 物理AI商业化节奏：富士通与发那科、安川、川崎的合作仍处探索阶段，无财务承诺；欧姆龙合作已进入实施阶段，直接受益于AI服务器PCB检测需求。
\end{itemize}
\subsubsection*{机构观点}
\paragraph{Bernstein} 数据中心项目取消与延迟数量在2026年前5个月已超2025年全年，预计仅60-65\%管线项目能按计划推进。电力设备交期仍紧张（燃气轮机排至2030年），但涨价幅度从30-50\%降至温和通胀。中国电气设备通过合作进入美国市场（西门子、GEV与XD合作）。冷却技术方面，CoolIT 15kW单相冷板在标准条件下验证可处理6倍于当前极限的热负荷，可能推迟两相DTC商业化窗口。
{\small\color{refgray}数据：2026年前5个月数据中心项目取消/延迟数量超过2025年全年\par}
{\small\color{refgray}数据：预计仅60-65\%管线项目能按计划推进\par}
{\small\color{refgray}数据：燃气轮机交期排至2030年，大型往复式发动机至2029-30年\par}
{\small\color{refgray}数据：涨价幅度从30-50\%降至温和通胀\par}
{\small\color{refgray}边际变化：从电力瓶颈转向项目落地瓶颈，公众反对成为主要障碍；单相DTC寿命预期延长，两相DTC商业化可能推迟。\par}
\paragraph{GS} 富士通与发那科、安川、川崎合作推进物理AI，计划2026年9月启动测试，AI服务器量产在2027年1-3月，FY3/28全面贡献盈利。欧姆龙合作已进入实施阶段，三项技术直接服务AI服务器PCB检测。UMT以17亿新台币收购DYNAHZ切入AI数据中心高速互联，DYNAHZ已进入AI客户体系，计划2028年前扩产2-3倍。
{\small\color{refgray}数据：富士通物理AI业务目标2035年新增3万亿日元收入\par}
{\small\color{refgray}数据：欧姆龙合作包含三项可直接部署的PCB检测技术\par}
{\small\color{refgray}数据：UMT收购价17亿新台币，DYNAHZ计划2028年扩产2-3倍\par}
{\small\color{refgray}数据：UMT 2025年营收24.52亿新台币，增速从47\%放缓至5\%\par}
{\small\color{refgray}边际变化：物理AI从平台开发进入系统集成与商业化阶段；欧姆龙合作已实现技术落地而非探索框架。\par}
\paragraph{MS} 2026年IT预算增速预期小幅回升至+3.8\%，1年期预算上修/下修比率自1Q24以来首次突破1.0x至1.2x。软件增速+4.1\%领跑，安全软件+12.4\%为最清晰赢家。微软在GenAI支出中份额最大但热度从32\%降至27\%。AI支出增量资金占比从61\%降至54\%，重新分配占比升至32\%。
{\small\color{refgray}数据：2026年IT预算增速预期+3.8\%，连续第三次调查改善\par}
{\small\color{refgray}数据：预算上修/下修比率从0.8x升至1.2x\par}
{\small\color{refgray}数据：软件增速+4.1\%，安全软件+12.4\%\par}
{\small\color{refgray}数据：AI支出增量资金占比从61\%降至54\%\par}
{\small\color{refgray}边际变化：预算上修比率转正，但AI支出从增量转向替代，资金流向更挑剔。\par}
\subsubsection*{关键数据}
\begin{itemize}
  \item 2026年前5个月数据中心项目取消/延迟数量超2025年全年
  \item 预计仅60-65\%管线项目能按计划推进
  \item CoolIT 15kW单相冷板验证可处理6倍于当前极限热负荷
  \item 2026年IT预算增速预期+3.8\%，上修/下修比率1.2x
  \item 安全软件增速+12.4\%，为最清晰赢家
  \item AI支出增量资金占比从61\%降至54\%
  \item 富士通物理AI目标2035年新增3万亿日元收入
  \item UMT收购DYNAHZ切入AI数据中心高速互联，计划2028年扩产2-3倍
\end{itemize}
\begin{figure}[htbp]
\centering
\includegraphics[width=0.88\linewidth]{figures/fig_011.jpg}
\caption{EXHIBIT 1 - EXHIBIT 2: Liquid cooling format by GPU TDP}
\end{figure}
\begin{figure}[htbp]
\centering
\includegraphics[width=0.88\linewidth]{figures/fig_145.jpg}
\caption{Exhibit 1 - Exhibit 1: External IT spending growth expectations over time \textbackslash{}- 2026 CIO IT budget growth expectations increased slightly to +3.8\% (up from 3.7\% in 1Q26). \textbackslash{}- Looking ahead into the next 3 years, c. 42\% of CIOs expected IT spendi}
\end{figure}
\begin{figure}[htbp]
\centering
\includegraphics[width=0.88\linewidth]{figures/fig_150.jpg}
\caption{Exhibit 1 - Exhibit 1: CIOs Indicate Expectations for Microsoft, LLM Providers, Amazon, and Salesforce to Gain the Largest Incremental Share of GenAI Spend in 2026... Vendors Gaining Largest Incremental Share of Gen AI Spend in 2026 Exhibit}
\end{figure}
\begin{figure}[htbp]
\centering
\includegraphics[width=0.88\linewidth]{figures/fig_151.jpg}
\caption{Exhibit 1 - Exhibit 3: At 3.8\%, CIOs Continue to Expect Modest Budget Growth Acceleration YoY (+14 bps from 3.7\% in 2025) Exhibit 4: Sequentially, Budget Growth Expectations Were Revised Up Across All Industries (Except Services)}
\end{figure}
{\small\color{refgray}本节综合 7 篇研究；涉及 Bernstein、GS、MS。}
\newpage
\subsection{创新药临床突破与医疗设备政策转向}
\textbf{中国创新药在ADC和联合疗法领域取得里程碑式进展，sac-TMT头对头击败Keytruda+化疗，Tango组合疗法在胰腺癌中实现92\%客观缓解率，同时中国医疗设备集采政策基调转向合理定价，日本医疗AI应用进入可量化效率阶段。}
\subsubsection*{主流共识}
\begin{itemize}
  \item ADC+PD-1联合疗法正在成为多个癌种的新标准治疗，sac-TMT在非小细胞肺癌中的III期成功验证了这一趋势。
  \item 中国医疗设备集采政策强调合理价格和全生命周期成本，而非单纯压低采购价，降价幅度预计比耗材温和。
  \item AI在药物发现和医疗流程中的应用已从概念验证进入实际效率提升阶段，但实验验证和高质量数据仍是核心壁垒。
\end{itemize}
\subsubsection*{分歧与条件差异}
\begin{itemize}
  \item 市场对sac-TMT全球商业潜力的定价存在分歧：MS预测41亿美元（55\%成功率），而Visible Alpha共识为59亿美元。
  \item 中国设备集采短期可能因各省准备期而延缓当前招标活动，但中长期将释放集中采购需求。
  \item 日本医疗AI补贴政策（最高8000万日元/院）能否转化为设备厂商的长期市场份额增长，而非同质化竞争。
  \item Tango的PRMT5抑制剂虽数据领先，但面临与Revolution Medicines合作的不确定性，已启动与ERAS的替代合作。
\end{itemize}
\subsubsection*{机构观点}
\paragraph{MS} 科伦博泰sac-TMT联合Keytruda在OptiTROP-Lung06 III期试验中达到PFS主要终点，覆盖PD-L1 TPS<1\%初治非鳞非小细胞肺癌患者，OS也观察到积极趋势。此前Lung05在PD-L1≥1\%人群中已成功，此次在更大患者群体中验证了联合疗法的潜力。全球临床布局加速，默沙东正推进17项全球III期研究。MS预测2035年风险调整后销售为41亿美元（55\%成功率），低于Visible Alpha共识的59亿美元，反映市场对多适应症成功概率的判断分歧。
{\small\color{refgray}数据：OptiTROP-Lung06：PFS主要终点达成，OS积极趋势\par}
{\small\color{refgray}数据：MS 2035年风险调整销售预测：41亿美元（成功率55\%）\par}
{\small\color{refgray}数据：Visible Alpha共识：59亿美元\par}
{\small\color{refgray}数据：默沙东推进17项全球III期研究\par}
{\small\color{refgray}边际变化：sac-TMT从PD-L1高表达人群扩展至低表达人群，临床验证范围扩大，全球里程碑关键拼图就位。\par}
\paragraph{GS} 中国大型医用设备集采通知要求各省2026年底前完成至少一轮乙类设备采购，节奏快于市场预期的2027-2028年。政策基调转向合理价格、质量适宜和全生命周期成本（含耗材、维保），而非单纯压低采购价。甲类设备由国家卫健委统一集采，乙类由省级推进。短期各省准备集采可能延缓当前招标活动，但国产龙头可通过份额提升和海外扩张部分对冲价格压力。
{\small\color{refgray}数据：2026年底前各省完成至少一轮乙类设备集采\par}
{\small\color{refgray}数据：甲类设备：MR-Linac、X射线立体定向放射手术系统\par}
{\small\color{refgray}数据：乙类设备：PET/MR、PET/CT、腔镜手术机器人\par}
{\small\color{refgray}数据：政策强调全生命周期成本而非单纯低价\par}
{\small\color{refgray}边际变化：集采时间表提前，但政策基调从唯低价转向合理定价，降价幅度预计比耗材温和。\par}
\paragraph{GS} 强生二季度营收253亿美元超预期，创新药164亿美元是主要驱动力。口服银屑病新药Icotyde上市90天内商业保险覆盖率超50\%，处方量超18000张，覆盖11000名患者，放量速度重新定义口服药竞争标准。这对日本武田同类口服药zasocitinib构成竞争，同时强生医疗科技板块有机增速3.7\%低于长期目标5-7\%，但管理层预计下半年回升。
{\small\color{refgray}数据：强生二季度营收253亿美元，创新药164亿美元\par}
{\small\color{refgray}数据：Icotyde：90天商业保险覆盖率>50\%\par}
{\small\color{refgray}数据：Icotyde处方量>18000张，覆盖>11000名患者\par}
{\small\color{refgray}数据：强生医疗科技有机增速3.7\%，低于5-7\%目标\par}
{\small\color{refgray}边际变化：Icotyde的快速放量表明口服药在皮肤科领域正替代生物制剂成为一线选择，支付端接纳速度超预期。\par}
\paragraph{JPM} 成都医疗调研覆盖9家公司，核心发现：中国创新驱动力从成本优势转向临床差异化。AI在药物发现中已成标配，但高质量物理数据和实验能力比算法更关键。对外授权交易活跃，跨国药企对中国资产的兴趣集中在临床差异化而非低成本。行业规范对合规公司影响有限，但将加速临床必需药与推广驱动型辅助药的分化。
{\small\color{refgray}数据：调研覆盖9家生物科技、CRO和器械公司\par}
{\small\color{refgray}数据：AI减少合成分子数量，但不取代实验验证\par}
{\small\color{refgray}数据：科伦博泰与默沙东、BioKin与BMS等授权案例\par}
{\small\color{refgray}数据：合规公司反馈：政策收紧对销售影响不大\par}
{\small\color{refgray}边际变化：中国医疗创新叙事从成本优势转向临床差异化，AI和授权交易成为新驱动力，行业规范加速分化。\par}
\paragraph{JPM} 科伦博泰sac-TMT联合Keytruda在OptiTROP-Lung06 III期试验中头对头击败Keytruda+化疗，这是ADC+PD-1组合首次在随机对照III期研究中证明优于当前标准治疗。覆盖PD-L1阴性非鳞非小细胞肺癌，患者群体大于此前Lung05的PD-L1高表达人群。中国NMPA已授予1L三阴性乳腺癌优先审评资格，数据将在ESMO和SABCS公布。JPM上调目标价至648港元，核心逻辑是临床不确定性下降。
{\small\color{refgray}数据：OptiTROP-Lung06：头对头击败Keytruda+化疗\par}
{\small\color{refgray}数据：覆盖PD-L1阴性非鳞非小细胞肺癌，人群更大\par}
{\small\color{refgray}数据：中国NMPA授予1L三阴性乳腺癌优先审评\par}
{\small\color{refgray}数据：JPM上调目标价至648港元（从535港元）\par}
{\small\color{refgray}边际变化：sac-TMT从有潜力的候选药物变为多个适应症中可验证的标准治疗替代方案，临床不确定性实质性下降。\par}
\paragraph{JPM} 日本医疗AI应用进入可量化效率阶段：一家医院通过AI生成电子病历将记录时间压缩至60秒以内，另一家使用Claude效率提升1.8倍，RPA自动化每月节省550工时（相当于3.6名全职员工）。2026年7月起日本开放最高8000万日元/院的ICT设备补贴，覆盖AI问诊、生成式AI文书等。但设备厂商需将AI功能转化为市场份额增长，避免同质化竞争。
{\small\color{refgray}数据：AI生成电子病历：记录时间<60秒\par}
{\small\color{refgray}数据：Claude辅助：效率提升1.8倍\par}
{\small\color{refgray}数据：RPA自动化：每月节省550工时，追回数千万日元漏报\par}
{\small\color{refgray}数据：ICT补贴：最高8000万日元/院\par}
{\small\color{refgray}边际变化：日本医院AI采纳从试点进入成本节约和收入回收阶段，政府补贴创造明确政策窗口。\par}
\paragraph{JPM} Tango Therapeutics的vopimetostat（PRMT5抑制剂）联合daraxonrasib（RAS抑制剂）在2L+胰腺导管腺癌患者中取得92\%客观缓解率和90\%6个月无进展生存率，远超daraxonrasib单药的30-35\%ORR和历史化疗的10-15\%。约22\%美国PDAC患者存在MTAP缺失，vopimetostat通过MTA协同机制选择性抑制肿瘤细胞。Tango已启动与ERAS的替代合作以对冲与Revolution Medicines的不确定性。
{\small\color{refgray}数据：组合疗法ORR：92\%\par}
{\small\color{refgray}数据：6个月PFS率：90\%\par}
{\small\color{refgray}数据：daraxonrasib单药ORR：30-35\%\par}
{\small\color{refgray}数据：历史化疗ORR：10-15\%，中位OS约6个月\par}
{\small\color{refgray}边际变化：PRMT5+RAS抑制剂组合将胰腺癌二线治疗疗效天花板大幅抬高，Tango通过多路径合作对冲竞争不确定性。\par}
\subsubsection*{关键数据}
\begin{itemize}
  \item OptiTROP-Lung06：PFS主要终点达成，OS积极趋势
  \item MS 2035年风险调整销售预测：41亿美元（成功率55\%），Visible Alpha共识59亿美元
  \item 中国设备集采：2026年底前各省完成至少一轮乙类采购
  \item Icotyde：90天商业保险覆盖率>50\%，处方量>18000张
  \item Tango组合疗法ORR：92\%，6个月PFS率90\%
  \item 日本AI应用：AI生成病历<60秒，RPA每月节省550工时
  \item ICT补贴：最高8000万日元/院
  \item JPM上调科伦博泰目标价至648港元
\end{itemize}
\begin{figure}[htbp]
\centering
\includegraphics[width=0.88\linewidth]{figures/fig_135.jpg}
\caption{Figure 1 - Figure 1: Clinical workflow support systems Japan market size Note: Historical data are Medley estimates. \#\# (1) AI significantly reduces documentation time in medical practices In one example at a hospital, use of AI enabled th}
\end{figure}
\begin{figure}[htbp]
\centering
\includegraphics[width=0.88\linewidth]{figures/fig_141.jpg}
\caption{Figure 6 - Figure 6: Vopimetostat (TNG462) relative to TNG908 \#\# A Short Primer on MTAP Deletion and PRMT5 As a Target in Cancer Methylthioadenosine phosphorylase (MTAP) is an enzyme encoded on chromosome 9p21, adjacent to the tumor suppre}
\end{figure}
\begin{figure}[htbp]
\centering
\includegraphics[width=0.88\linewidth]{figures/fig_063.jpg}
\caption{Exhibit 1 - Exhibit 2: Dental implant renewal VBP started in July; expecting additional announcements towards year end 2026}
\end{figure}
\begin{figure}[htbp]
\centering
\includegraphics[width=0.88\linewidth]{figures/fig_064.jpg}
\caption{Exhibit 1 - Exhibit 2: Dental implant renewal VBP started in July; expecting additional announcements towards year end 2026}
\end{figure}
{\small\color{refgray}本节综合 7 篇研究；涉及 MS、GS、JPM。}
\newpage
\subsection{工业与新能源：分化加剧，结构性机会浮现}
\textbf{全球工业与新能源板块正经历显著分化：中国汽车市场受补贴透支影响持续调整，但出口与新能源渗透率结构性提升；光伏产业链价格阴跌，仅逆变器与薄膜环节维持盈利；铜矿运营协同、低轨卫星、AI数据中心等细分领域则展现出独立于宏观的增量机会。}
\subsubsection*{主流共识}
\begin{itemize}
  \item 中国汽车市场国内零售受补贴透支影响，短期难以恢复至正常年化2200万辆水平，出口成为主要增长引擎。
  \item 光伏产业链供需失衡持续，多晶硅、硅片、组件价格仍在下行，库存累积未见缓解。
  \item 全球铜矿供应紧张背景下，资产层面的运营整合（如BHP智利铜矿合作）可释放显著协同价值。
\end{itemize}
\subsubsection*{分歧与条件差异}
\begin{itemize}
  \item 中国汽车市场：部分机构认为政策支持边际效应递减，行业需自然修复；另一些则关注新车型周期（如小鹏MONA系列）能否驱动销量加速回升。
  \item 光伏行业：Citi强调逆变器与薄膜环节是仅有的盈利点，而主产业链企业仍深度亏损；但GS对低轨卫星射频组件等新兴领域持乐观态度，认为其将重塑相关公司估值。
  \item 工业自动化：MS对埃斯顿自动化持谨慎态度，认为扣非利润虽转正但估值偏高；而JPM对华勤技术的数据中心业务转型更为积极，预计2026下半年起利润贡献加速。
\end{itemize}
\subsubsection*{机构观点}
\paragraph{Bernstein} 中国汽车市场6月零售同比下滑21.4\%，年化销量2070万辆，低于正常需求2200万辆。补贴政策提前释放超500万辆需求，叠加高基数效应，短期缺乏回升催化剂。出口同比增长80.2\%，成为相对亮点。新能源渗透率达59.2\%，但插电混动销量同比下滑28.7\%，因纯电续航提升、价差缩小及政策优惠调整。
{\small\color{refgray}数据：6月零售171万辆，同比-21.4\%\par}
{\small\color{refgray}数据：年化零售2070万辆，低于正常2200万辆\par}
{\small\color{refgray}数据：补贴提前释放超500万辆需求\par}
{\small\color{refgray}数据：出口同比+80.2\%\par}
{\small\color{refgray}边际变化：从关注国内零售下滑转向强调出口增长和新能源内部结构分化（纯电优于插混）。\par}
\paragraph{Bernstein} 印度两轮车电动化不会改变市场底层竞争逻辑。大众市场（通勤摩托/家庭踏板）仍由规模、成本、服务网络决定胜负，Hero等传统龙头凭借自我强化循环占据主导。品牌专属品类（如Bajaj Pulsar、本田Activa）则依赖生态护城河。电动化可能使大众市场更分散，但结构性优势仍属于具备规模与成本控制能力的玩家。
{\small\color{refgray}数据：印度两轮车年销超2000万辆，前十品牌占2/3份额\par}
{\small\color{refgray}数据：前三品牌占约40\%份额\par}
{\small\color{refgray}数据：Hero在通勤摩托市场统治地位数十年\par}
{\small\color{refgray}数据：本田用约15年才在通勤市场达约1/4份额\par}
{\small\color{refgray}边际变化：从电动化颠覆论转向强调传统竞争要素（规模、成本、渠道）的持续性。\par}
\paragraph{MS} 2026年7月起社保征缴移交税务部门，长三角汽车零部件企业首当其冲。按实际工资缴纳社保将导致雇主劳动力成本上升约25\%，员工到手收入减少约10\%。经销商（净利率1-2\%）和劳动密集型零部件商最脆弱，芯片、软件等环节相对免疫。双重冲击企业利润与家庭购买力，市场低估了消费端影响。
{\small\color{refgray}数据：上海社保缴费下限7500元/月，江苏约4800元/月\par}
{\small\color{refgray}数据：雇主劳动力成本上升约25\%\par}
{\small\color{refgray}数据：员工到手收入减少约10\%\par}
{\small\color{refgray}数据：经销商净利率仅1-2\%\par}
{\small\color{refgray}边际变化：从关注企业端成本转向同时关注居民可支配收入对汽车消费的间接压制。\par}
\paragraph{MS} 埃斯顿自动化2Q26净利润中位数约6700万元，高于市场预期但低于MS预测。扣非净利润约4800万元，同比扭亏、环比+149\%，但盈利规模相对市值仍偏薄。增长来自产品结构改善及一次性公允价值收益。人形机器人业务单独估值（15倍2028年PS），但规模小、份额低，估值低于行业均值。
{\small\color{refgray}数据：2Q26净利润中位数6700万元\par}
{\small\color{refgray}数据：扣非净利润4800万元，环比+149\%\par}
{\small\color{refgray}数据：核心业务估值50倍PE，低于历史均值65倍\par}
{\small\color{refgray}数据：人形机器人业务按2028年营收6.8亿、15倍PS估值\par}
{\small\color{refgray}边际变化：从关注报表盈利转向强调扣非利润的可持续性与估值合理性。\par}
\paragraph{MS} 思源电气1H26收入+27\%、净利润+15\%，低于预期，主因汇率与少数股东权益变动。但北美订单突破是核心亮点：全球变压器短缺下，海外订单增速50-100\%，预计海外变压器订单从2025年约40亿增至2026/27年60亿/100亿。产能扩张至4Q27，新产品（1000kV GIS、特高压直流、超级电容）打开第二曲线。
{\small\color{refgray}数据：1H26国内订单同比+25\%\par}
{\small\color{refgray}数据：海外订单增速50-100\%\par}
{\small\color{refgray}数据：海外变压器订单2025年约40亿，2026年超60亿，2027年有望破100亿\par}
{\small\color{refgray}数据：产能扩张至4Q27\par}
{\small\color{refgray}边际变化：从关注短期汇兑损益转向北美订单突破与产品结构升级。\par}
\paragraph{Citi} 中国光伏产业链价格持续阴跌：n型多晶硅周跌0.8\%至31.6元/kg，硅片跌1.1\%，组件跌0.6\%。供需失衡未解，7月多晶硅产量预计超10万吨，月需求仅9.1万吨，库存达31万吨。上半年业绩预告显示，仅逆变器（德业股份净利润同比+75-79\%）和光伏薄膜（福斯特净利润同比+75.4\%）维持盈利，主产业链企业深度亏损。
{\small\color{refgray}数据：n型多晶硅31.6元/kg，周跌0.8\%\par}
{\small\color{refgray}数据：182mm硅片0.86元/W，周跌1.1\%\par}
{\small\color{refgray}数据：组件0.71元/W，周跌0.6\%\par}
{\small\color{refgray}数据：7月多晶硅产量预计超10万吨，需求9.1万吨\par}
{\small\color{refgray}边际变化：从关注全行业亏损转向识别逆变器与薄膜环节的结构性盈利。\par}
\paragraph{GS} BHP旗下Spence铜矿与Sierra Gorda合资项目（KGHM 55\%/South32 45\%）深度协作可释放20-40亿美元协同价值，相当于各自NPV的10-20\%。核心来源：Sierra Gorda氧化矿堆存（含32万吨铜）利用Spence闲置阴极产能（年13万吨），每年可增约2亿美元EBITDA；尾矿处理可推迟超5亿美元资本支出；采购与物流年省约4000万美元。
{\small\color{refgray}数据：协同价值20-40亿美元\par}
{\small\color{refgray}数据：氧化矿堆存含32万吨铜\par}
{\small\color{refgray}数据：Spence闲置阴极产能13万吨/年\par}
{\small\color{refgray}数据：每年EBITDA增量约2亿美元\par}
{\small\color{refgray}边际变化：从关注大型并购转向资产层面运营整合的低资本支出协同模式。\par}
\paragraph{GS} 通宇通讯2Q26净亏损1200万元，较1Q26的1300万元收窄，但低于GS预期的1100万元盈利。公司正从基站射频转向低轨卫星射频组件，GS上调全球LEO卫星安装基数预测：2028年2.4万台、2031年30.5万台，中国占比从2025年3\%升至2031年8\%。卫星组件收入占比提升将驱动毛利率U型反转，2026年毛利率下调只是短期扰动。
{\small\color{refgray}数据：2Q26净亏损1200万元\par}
{\small\color{refgray}数据：全球LEO卫星安装基数2028年2.4万台、2031年30.5万台\par}
{\small\color{refgray}数据：中国安装基数从2025年253台增至2031年2.375万台\par}
{\small\color{refgray}数据：卫星网络设备收入从2025年5400万元增至2028年26.4亿元\par}
{\small\color{refgray}边际变化：从关注短期亏损收窄转向低轨卫星星座建设带来的长期需求曲线。\par}
\paragraph{GS} 小鹏MONA L03 SUV起售价12.38万元，开放7分钟获2万不可退订单，1小时达4.7万，订单速度是M03两倍以上。M03在10-15万元纯电轿车市场月均销1.23万辆排名第一。GS预计新车型周期将推动季度销量同比增速从个位数回升至28-59\%。L03增程版在10-15万元SUV市场差异化竞争，EREV版本竞争力突出。
{\small\color{refgray}数据：L03起售价12.38万元，比预售价低2万元\par}
{\small\color{refgray}数据：7分钟2万不可退订单，1小时4.7万\par}
{\small\color{refgray}数据：M03月均销1.23万辆，10-15万纯电轿车市场第一\par}
{\small\color{refgray}数据：预计季度销量同比增速28-59\%\par}
{\small\color{refgray}边际变化：从关注行业整体调整转向新车型周期驱动的销量加速与差异化定价策略。\par}
\paragraph{JPM} 华勤技术2Q26净利润中值19亿元，同比+83\%、环比+81\%，超预期65\%，但扣非净利润8.6亿元（同比+15\%、环比+5\%）。市场担忧智能手机需求变化，但PC、可穿戴、AIoT份额提升抵消影响。核心看点：数据中心业务从2026下半年起加速，机柜级AI项目量产推动收入同比+50\%以上，预计2027/28年贡献毛利的33\%/38\%。
{\small\color{refgray}数据：2Q26净利润19亿元，超预期65\%\par}
{\small\color{refgray}数据：扣非净利润8.6亿元，同比+15\%\par}
{\small\color{refgray}数据：数据中心收入2026年同比+50\%以上\par}
{\small\color{refgray}数据：数据中心2027/28年贡献毛利33\%/38\%\par}
{\small\color{refgray}边际变化：从关注消费电子需求波动转向数据中心业务作为利润引擎的结构性转型。\par}
\subsubsection*{关键数据}
\begin{itemize}
  \item 中国汽车6月零售同比-21.4\%，年化2070万辆，低于正常2200万辆
  \item 补贴提前释放超500万辆需求
  \item 光伏多晶硅库存31万吨，周环比+3\%
  \item 逆变器环节德业股份1H26净利润同比+75-79\%
  \item BHP智利铜矿协同价值20-40亿美元
  \item 全球LEO卫星安装基数2031年达30.5万台
  \item 小鹏L03 1小时获4.7万不可退订单
  \item 华勤数据中心2027/28年贡献毛利33\%/38\%
\end{itemize}
\begin{figure}[htbp]
\centering
\includegraphics[width=0.88\linewidth]{figures/fig_006.jpg}
\caption{Exhibit 2 - Exhibit 3 - Exhibit 4) EXHIBIT 2: Chinese autos retail sales came in at 1.71mn units in June 2026, -21.4\% yoy EXHIBIT 3: BYD Group (incl. BYD, Denza, Fangchengbao, Yangwang, and Linghui) came back to the top with 224k units, 13.2\%}
\end{figure}
\begin{figure}[htbp]
\centering
\includegraphics[width=0.88\linewidth]{figures/fig_058.jpg}
\caption{Figure 1 - \textbackslash{}- JA Solar estimates its net loss to be Rmb 2.4-2.9 bn in 1H26E, including Rmb1.3-1.8bn in 2Q26E (vs. losses of Rmb942m in 2Q25 and Rmb1,067m in 1Q26). Figure 1. ePRC weekly polysilicon prices in the week ended Jul 15 Figure 2. PRC weekly polysilicon price © }
\end{figure}
\begin{figure}[htbp]
\centering
\includegraphics[width=0.88\linewidth]{figures/fig_070.jpg}
\caption{Exhibit 1 - Exhibit 1: BHP's Spence copper mine and the Sierra Gorda (KGHM 55\%/South32 45\%) in Northern Chile are just 20km apart Exhibit 2: Key potential synergies could reach US\textbackslash{}\$2-4bn: Our assessment of the two mines indicates there could}
\end{figure}
\begin{figure}[htbp]
\centering
\includegraphics[width=0.88\linewidth]{figures/fig_104.jpg}
\caption{Exhibit 1 - Exhibit 1: Tongyu's waterfall chart for 2025-30E Satellite networking equipment as key drivers ahead Exhibit 2: Tongyu's revenue and GM trends Global LEO satellite TAM: We updated our Global LEO satellite TAM in Jul (report lin}
\end{figure}
{\small\color{refgray}本节综合 10 篇研究；涉及 Bernstein、MS、Citi、GS、JPM。}
\newpage
\subsection{金融与资金流向：信用结构分化与定价权转移}
\textbf{6月中国金融数据显示信用总量放缓但结构分化加剧：住户部门持续去杠杆，企业贷款“去量求质”，银行定价权被动削弱，净息差承压。保险板块资金流向亦呈分歧，南向资金偏好高股息标的，长线资金则更关注基本面风险。}
\subsubsection*{主流共识}
\begin{itemize}
  \item 6月社融与人民币贷款增速均低于预期，信用扩张放缓。
  \item 住户部门仍在去杠杆，上半年住户贷款净减少3670亿元，居民加杠杆意愿低迷。
  \item 新发放企业贷款利率降至约3.0\%，银行定价能力持续走弱。
\end{itemize}
\subsubsection*{分歧与条件差异}
\begin{itemize}
  \item 贷款增速放缓是政策主动“去量求质”的积极信号，还是需求疲弱的负面反映？
  \item 银行净息差走向取决于贷款定价下行与存款成本下降的速度对比，方向存在不确定性。
  \item 保险板块中，南向资金与长线资金对同一标的持仓方向相反，反映不同资金属性对估值与基本面的定价逻辑差异。
\end{itemize}
\subsubsection*{机构观点}
\paragraph{MS} 6月社融增速降至7.4\%，人民币贷款增速降至5.3\%，但这是政策主动优化结构、注重质量的延续。住户贷款环比回升但同比仍降56\%，上半年净减少3670亿元，居民去杠杆节奏从骤降转为缓慢调整。企业中长期贷款上半年同比少增23\%，反映银行在取消小微贷款增长目标后的理性放贷。贷款增速放缓反而有利于银行资产回报表现。新发放企业贷款利率约3\%，按揭利率约3.1\%，资金价格下行但存款成本也在下降，净息差走向取决于两端调整速度。
{\small\color{refgray}数据：6月社融增速7.4\%（前值7.7\%）\par}
{\small\color{refgray}数据：人民币贷款增速5.3\%（前值5.5\%）\par}
{\small\color{refgray}数据：上半年住户贷款净减少3670亿元\par}
{\small\color{refgray}数据：企业中长期贷款上半年同比少增23\%\par}
{\small\color{refgray}边际变化：此前市场担忧信用收缩，但MS强调这是政策主动“去量求质”，且贷款增速放缓可能改善银行资产回报，边际上从悲观转向中性偏积极。\par}
\paragraph{JPM} 6月社融和新增人民币贷款双双低于预期，M1-M2剪刀差走阔至-4.0\%。但核心问题不在总量，而在于银行正在失去对贷款定价的控制力。新发放企业贷款加权平均利率降至约3.0\%，同比下降约20bp，且发生在没有进一步降息的环境下，表明银行主动或被动让利。零售贷款余额同比收缩1.3\%，短期零售贷款增速降至-6.7\%，中长期零售贷款增速放缓至0.5\%。高收益资产收缩叠加低收益资产占比上升，正在挤压净息差。预计二季度净息差可能环比持平或小幅下降，逆转一季度环比改善3bp的趋势。
{\small\color{refgray}数据：6月新发放企业贷款利率约3.0\%，同比降20bp\par}
{\small\color{refgray}数据：零售贷款余额同比收缩1.3\%\par}
{\small\color{refgray}数据：短期零售贷款增速-6.7\%\par}
{\small\color{refgray}数据：中长期零售贷款增速0.5\%\par}
{\small\color{refgray}边际变化：此前市场关注降息对息差的直接影响，JPM揭示即使政策利率不变，贷款结构变化（零售收缩、对公占比上升）本身就在压缩息差，边际上更强调定价权而非利率水平。\par}
\paragraph{MS} 6月保险板块资金流向呈现清晰分化。南向资金持续流入中国财险（持股比例升1.1ppt），中国太平亦增0.3ppt，但中国人寿、平安、太保南向持股各降0.7ppt。长线资金方面，平安持仓环比降0.6ppt至2.8\%，友邦降0.3ppt至2.6\%，但友邦南向资金在6月反弹0.2ppt。中国财险约6\%的股息率是吸引南向资金的核心因素，尽管存在巨灾损失担忧。长线资金更关注跨境业务（MCV）风险和AI板块轮动影响。
{\small\color{refgray}数据：中国财险南向持股比例6月升1.1ppt\par}
{\small\color{refgray}数据：中国平安长线持仓环比降0.6ppt至2.8\%\par}
{\small\color{refgray}数据：友邦保险长线持仓降0.3ppt至2.6\%\par}
{\small\color{refgray}数据：友邦南向资金6月反弹0.2ppt\par}
{\small\color{refgray}边际变化：此前市场关注保险板块整体资金流向，MS指出南向与长线资金在同一标的上方向相反，边际上强调不同资金属性的定价逻辑差异，而非单一趋势。\par}
\subsubsection*{关键数据}
\begin{itemize}
  \item 6月社融增速7.4\%（前值7.7\%）
  \item 人民币贷款增速5.3\%（前值5.5\%）
  \item 上半年住户贷款净减少3670亿元
  \item 新发放企业贷款利率约3.0\%，同比降20bp
  \item M1-M2剪刀差-4.0\%（前值-3.1\%）
  \item 零售贷款余额同比收缩1.3\%
  \item 中国财险南向持股比例6月升1.1ppt
  \item 中国平安长线持仓环比降0.6ppt至2.8\%
\end{itemize}
\begin{figure}[htbp]
\centering
\includegraphics[width=0.88\linewidth]{figures/fig_042.jpg}
\caption{Exhibit 1 - Exhibit 1: Southbound holdings of H-share listed insurers and changes in June 2026 Exhibit 2: H-share listed insurers' holdings among long-only EM and China active managers}
\end{figure}
\begin{figure}[htbp]
\centering
\includegraphics[width=0.88\linewidth]{figures/fig_125.jpg}
\caption{Figure 1 - Figure 3: Corporate and government bond financing growth edged down to 12.8\% y/y in Jun-26 vs 13.3\% y/y in May-26}
\end{figure}
\begin{figure}[htbp]
\centering
\includegraphics[width=0.88\linewidth]{figures/fig_158.jpg}
\caption{Exhibit 2 - Exhibit 2: Household deposit growth slowed to 7.1\% yoy in June. Both M1 and M2 growth moderated, to 4\% and 8\% yoy, respectively \#\# Disclosure Section}
\end{figure}
\begin{figure}[htbp]
\centering
\includegraphics[width=0.88\linewidth]{figures/fig_043.jpg}
\caption{Exhibit 1 - Exhibit 1: Southbound holdings of H-share listed insurers and changes in June 2026 Exhibit 2: H-share listed insurers' holdings among long-only EM and China active managers \#\# Disclosure Section}
\end{figure}
{\small\color{refgray}本节综合 3 篇研究；涉及 MS、JPM。}
\newpage
\subsection{IT服务与企业软件：AI资金再分配重塑增长格局}
\textbf{全球IT服务预算增速放缓，AI支出正从纯增量转向预算内再分配，企业软件领域AI变现闭环加速形成，但芯片供给和客户集中度构成结构性约束。}
\subsubsection*{主流共识}
\begin{itemize}
  \item 2026年全球IT预算增速小幅改善至3.8\%，但IT服务支出增速预期降至1.8\%，显示资金正从传统服务向软件和AI迁移。
  \item AI支出仍具增量性，但资金来源正从净新增预算转向现有预算再分配，软件预算成为主要腾挪来源。
  \item 企业软件领域AI变现加速，Zoom客户体验业务（ZCX）ARR突破2亿美元且增速加快，AI语音功能成为差异化壁垒。
\end{itemize}
\subsubsection*{分歧与条件差异}
\begin{itemize}
  \item AI资金再分配对IT服务商的影响：MS认为存量争夺加剧，传统外包面临被削减风险；GS则认为资本开支是营收增长领先指标，激进capex可支撑更高增速。
  \item IT服务增速放缓的持续性：MS调查显示全球IT服务支出预期环比下降20bp，但GS对金山云的测算显示AI云营收下半年有望重新加速至71\%，节奏差异源于基数效应和芯片供给。
  \item 企业软件AI变现的定价模式：JPM观察到Zoom从固定席位向消费量/结果导向迁移，而MS调查中仅33\%的CIO有净新增AI预算，显示定价模式转型仍面临预算约束。
\end{itemize}
\subsubsection*{机构观点}
\paragraph{MS} 2026年全球IT服务预算增速预期降至1.8\%，较上季下降20bp，而整体IT预算增速改善至3.8\%。AI资金从净新增预算（33\%，低于上季37\%）转向预算再分配（32\%），其中15\%来自软件预算。GenAI部署以IT运维（47\%）和软件开发（44\%）领先。公共云负载预计从当前48\%升至2028年底66\%。供应商折扣压力持续，35\%的CIO认为折扣力度加大。
{\small\color{refgray}数据：2026年IT服务支出增长预期1.8\%，环比下降20bp\par}
{\small\color{refgray}数据：整体IT预算增速预期3.8\%，接近历史均值4.1\%\par}
{\small\color{refgray}数据：仅33\%的CIO表示AI资金来自净新增预算，低于上季37\%\par}
{\small\color{refgray}数据：32\%的CIO通过预算再分配支持AI，其中15\%来自软件预算\par}
{\small\color{refgray}边际变化：AI融资方式从净新增预算主导转向预算再分配主导，IT服务支出预期增速首次出现环比下降，表明AI对传统服务的替代效应开始显现。\par}
\paragraph{GS} 金山云二季度营收预计同比增长27\%至30亿元，较一季度37\%放缓，主要受芯片供给偏紧和基数效应影响。AI云营收预计同比增长62\%，低于一季度的90\%，但下半年有望重新加速至71\%。管理层给出的150-200亿元全年capex指引可支撑约35\%的营收增长。小米+金山生态收入占比预计从2025年27\%升至2028年40\%。当前PS仅1.5x（2026E），低于历史均值2.5x两个标准差。
{\small\color{refgray}数据：Q2营收预计同比+27\%至30亿元，较Q1的+37\%放缓\par}
{\small\color{refgray}数据：AI云营收预计同比+62\%，低于Q1的+90\%\par}
{\small\color{refgray}数据：下半年AI云营收增速有望重新加速至71\%\par}
{\small\color{refgray}数据：全年capex指引150-200亿元，可支撑约35\%营收增长\par}
{\small\color{refgray}边际变化：市场将capex指引视为成本不确定性，GS将其重新解读为营收增长领先指标，并指出芯片供给而非需求是增速节奏变化的主因。\par}
\paragraph{JPM} Zoom客户体验业务（ZCX）ARR已突破2亿美元，过去两个季度增速持续加快至高双位数高端，驱动因素为AI语音产品ZVA的加入。近十个大单中九个为竞争性替换，九个附带付费AI功能。Zoom Phone ARR约7.5-8亿美元，年增长10-15\%，市场仍有2.8亿席位待转化，其中一半仍在使用传统电话系统。定价模式从固定席位向消费量/结果导向迁移。
{\small\color{refgray}数据：ZCX ARR突破2亿美元，增速加快至高双位数高端\par}
{\small\color{refgray}数据：近十个大单中九个为竞争性替换，九个附带付费AI功能\par}
{\small\color{refgray}数据：Zoom Phone ARR约7.5-8亿美元，年增长10-15\%\par}
{\small\color{refgray}数据：市场仍有2.8亿席位待转化，一半仍在使用传统电话系统\par}
{\small\color{refgray}边际变化：ZCX增速从五年前跨过1亿美元ARR后持续加速，AI语音功能ZVA成为差异化壁垒，定价模式向消费量/结果导向迁移，表明AI变现闭环已形成。\par}
\subsubsection*{关键数据}
\begin{itemize}
  \item 2026年IT服务支出增长预期1.8\%，环比下降20bp
  \item 仅33\%的CIO表示AI资金来自净新增预算，低于上季37\%
  \item 金山云AI云营收Q2同比+62\%，下半年有望加速至71\%
  \item Zoom ZCX ARR突破2亿美元，增速加快至高双位数高端
  \item Zoom Phone ARR约7.5-8亿美元，市场仍有2.8亿席位待转化
  \item 金山云全年capex指引150-200亿元，可支撑约35\%营收增长
  \item 小米+金山生态收入占比预计从27\%升至40\%
\end{itemize}
\begin{figure}[htbp]
\centering
\includegraphics[width=0.88\linewidth]{figures/fig_048.jpg}
\caption{Exhibit 1 - Exhibit 2: Sequentially, Budget Growth Expectations Rose Across All Sectors (Except Services) 2026 External IT Spending Growth Expectations by Sector}
\end{figure}
\begin{figure}[htbp]
\centering
\includegraphics[width=0.88\linewidth]{figures/fig_091.jpg}
\caption{Exhibit 1 - Exhibit 1: We model KC's revenue growth at \$27\textbackslash{}\%\$ yoy in 2Q26E moderating from \$37\textbackslash{}\%\$ yoy in 1Q26 due to chip constraints Exhibit 2: We expect AI related revenue contribution to increase to 40\%+ in 2026E vs. 31\% in 2025}
\end{figure}
\begin{figure}[htbp]
\centering
\includegraphics[width=0.88\linewidth]{figures/fig_049.jpg}
\caption{Exhibit 2 - Exhibit 2: Sequentially, Budget Growth Expectations Rose Across All Sectors (Except Services) 2026 External IT Spending Growth Expectations by Sector Exhibit 3: CIOs Continue to Expect Software to Be the Fastest Growing IT Segmen}
\end{figure}
\begin{figure}[htbp]
\centering
\includegraphics[width=0.88\linewidth]{figures/fig_050.jpg}
\caption{Exhibit 3 - Exhibit 5: CIOs in the Business Services, Financials, and Healthcare Industries Expect the Highest IT Budget Growth in vs 2025: External IT Spending by Vertical}
\end{figure}
{\small\color{refgray}本节综合 3 篇研究；涉及 MS、GS、JPM。}
\newpage
\subsection{印度即时零售：从跑马圈地到存量博弈}
\textbf{印度即时零售行业年化GMV已达约140亿美元，但高速扩张阶段已结束，行业正从增量竞争转向存量竞争。一线城市暗店密度已超盈利性均衡点约20\%，新增门店中90\%开在已有覆盖区域，单店订单密度面临摊薄压力。未来增长将依赖存量用户消费提升，而非新增用户，行业整合与盈利改善成为关键。}
\subsubsection*{主流共识}
\begin{itemize}
  \item 印度即时零售行业年化GMV约140亿美元，高速增长阶段已基本结束，未来增长主要依赖存量用户消费提升。
  \item 一线城市暗店密度已超过盈利性均衡点约20\%，新增门店中90\%开在已有覆盖区域，竞争加剧导致单店订单密度承压。
  \item 行业整体EBITDA利润率约-10\textasciitilde{}12\%，年化亏损约25亿美元，持续烧钱模式难以维持，补贴可能逐步减少。
  \item 到2030财年，可盈利市场规模约350亿美元，但需要4年时间实现，行业格局基本形成。
\end{itemize}
\subsubsection*{分歧与条件差异}
\begin{itemize}
  \item Blinkit已实现盈亏平衡，而其余玩家亏损较深，行业内部盈利能力分化明显。
  \item Blinkit75\%新店开在非一线城市，测试下沉市场；Zepto和亚马逊继续深耕高密度城市，单城20家店；Instamart和Blinkit单城约9家店，走广度路线。
  \item 一线城市同时被五大玩家覆盖的邮政编码占比从26\%跃升至44\%，仅由单一玩家服务的邮政编码已降至9\%，竞争格局快速变化。
\end{itemize}
\subsubsection*{机构观点}
\paragraph{Bernstein} 印度即时零售行业年化GMV约140亿美元，但高速增长阶段已结束。基于约1.9万个邮政编码的专有数据库，供给端推算盈利性可寻址市场在2030财年约350亿美元，行业已走到40\%。过去一季度新增约900家暗店，90\%开在已有覆盖区域，一线城市暗店密度超盈利性均衡点约20\%。未来增长依赖存量用户消费提升，行业整合与盈利改善成为关键。
{\small\color{refgray}数据：印度即时零售行业年化GMV约140亿美元\par}
{\small\color{refgray}数据：盈利性可寻址市场在2030财年约350亿美元，行业已走到40\%\par}
{\small\color{refgray}数据：过去一季度新增约900家暗店，90\%开在已有覆盖区域\par}
{\small\color{refgray}数据：一线城市暗店密度超盈利性均衡点约20\%\par}
{\small\color{refgray}边际变化：此前市场关注即时零售的快速增长和暗店扩张，但Bernstein指出行业已从增量竞争转向存量竞争，一线城市暗店密度超盈利性均衡点，新增门店主要摊薄现有订单密度，而非开拓新市场。\par}
\subsubsection*{关键数据}
\begin{itemize}
  \item 印度即时零售行业年化GMV约140亿美元
  \item 盈利性可寻址市场在2030财年约350亿美元，行业已走到40\%
  \item 过去一季度新增约900家暗店，90\%开在已有覆盖区域
  \item 一线城市暗店密度超盈利性均衡点约20\%
  \item 一线城市同时被五大玩家覆盖的邮政编码占比从26\%跃升至44\%
  \item 仅由单一玩家服务的邮政编码已降至9\%
  \item 行业整体EBITDA利润率约-10\textasciitilde{}12\%
  \item 年化亏损约25亿美元
\end{itemize}
\begin{figure}[htbp]
\centering
\includegraphics[width=0.88\linewidth]{figures/fig_016.jpg}
\caption{Exhibit 6 - EXHIBIT 2: Our methodology for QC SAM estimation using Supply Side economics of dark stores}
\end{figure}
\begin{figure}[htbp]
\centering
\includegraphics[width=0.88\linewidth]{figures/fig_017.jpg}
\caption{EXHIBIT 1 - EXHIBIT 2: Our methodology for QC SAM estimation using Supply Side economics of dark stores \textbackslash{}\# Eligible Pin Codes by tier and corresponding \# Dark Stores}
\end{figure}
\begin{figure}[htbp]
\centering
\includegraphics[width=0.88\linewidth]{figures/fig_018.jpg}
\caption{EXHIBIT 2 - EXHIBIT 2: Our methodology for QC SAM estimation using Supply Side economics of dark stores \textbackslash{}\# Eligible Pin Codes by tier and corresponding \# Dark Stores EXHIBIT 4: NOV growth trajectory has been reducing gradually across pla}
\end{figure}
\begin{figure}[htbp]
\centering
\includegraphics[width=0.88\linewidth]{figures/fig_019.jpg}
\caption{EXHIBIT 3 - EXHIBIT 4: NOV growth trajectory has been reducing gradually across players EXHIBIT 5: QoQ growth has been volatile across players with Zepto clearly pulling ahead in the last 2 quarters}
\end{figure}
{\small\color{refgray}本节综合 1 篇研究；涉及 Bernstein。}
\newpage
\subsection{AI基础设施去杠杆：健康重置而非泡沫破裂}
\textbf{AI板块近期调整并非趋势逆转，而是资产负债表驱动的健康去杠杆过程；硬件瓶颈持续至2028年，定价权仍在供给侧。}
\subsubsection*{主流共识}
\begin{itemize}
  \item 全球超大规模云厂商资产负债率约40\%，远低于历史泡沫峰值（中国地产100\%、美国电信230\%），AI基础设施杠杆水平安全。
  \item 大语言模型能力持续提升，中国头部模型正追赶美国同行，每代新模型解锁增量应用场景，需求创造仍是上行期权。
  \item 关键AI芯片产能紧张至少持续至2027-2028年，华为昇腾960预计2027年Q4推出，CXMT HBM3E同年量产，全球HBM供应拐点在2028年。
  \item A股IT板块融资交易占比从12\%降至8-9\%，杠杆资金出清接近尾声，ETF资金开始回流。
\end{itemize}
\subsubsection*{分歧与条件差异}
\begin{itemize}
  \item 市场部分观点认为AI调整是泡沫破裂前兆，但JPM认为回调是健康的去杠杆过程，而非趋势转变。
  \item 关于需求饱和时点存在分歧：JPM认为需求创造仍为上行期权，而部分投资者担忧企业级应用落地慢于预期。
  \item 对硬件瓶颈缓解节奏看法不一：JPM强调2028年前不会缓解，但部分机构认为台积电亚利桑那Fab 2在2027年进入3nm量产可提前缓解。
\end{itemize}
\subsubsection*{机构观点}
\paragraph{JPM} AI板块调整是资产负债表驱动的健康去杠杆，而非泡沫破裂。中国及全球云厂商负债率约40\%，远低于中国地产（100\%）和美国电信（230\%）峰值。模型能力持续提升，中国头部模型正追赶美国同行，每代新模型解锁增量应用场景。硬件瓶颈至少持续至2027-2028年：华为昇腾960预计2027年Q4推出，CXMT HBM3E同年量产，全球HBM供应拐点在2028年。A股IT板块融资交易占比从12\%降至8-9\%，杠杆出清接近尾声，ETF资金开始回流。未来两个季度不是验证AI资本支出可持续性的窗口，定价权仍在供给侧。
{\small\color{refgray}数据：全球超大规模云厂商资产负债率约40\%\par}
{\small\color{refgray}数据：中国房地产企业2021年峰值负债率接近100\%\par}
{\small\color{refgray}数据：美国电信公司2001年互联网泡沫时负债率高达230\%\par}
{\small\color{refgray}数据：华为昇腾960预计2027年Q4推出\par}
{\small\color{refgray}边际变化：市场此前担忧AI调整是泡沫破裂前兆，JPM反转为健康去杠杆叙事，强调资产负债表稳健和硬件瓶颈持续，将关注点从需求端转向供给侧约束。\par}
\subsubsection*{关键数据}
\begin{itemize}
  \item 全球超大规模云厂商资产负债率约40\%，远低于中国地产（100\%）和美国电信（230\%）峰值
  \item 华为昇腾960预计2027年Q4推出，CXMT HBM3E 12层堆叠2027年量产
  \item 全球HBM供应大规模释放拐点预计2028年
  \item A股IT板块融资交易占比从12\%降至8-9\%，杠杆出清接近尾声
  \item A股换手率从6.5\%高点回落但未达恐慌水平
  \item 台积电亚利桑那Fab 2在2027年进入3nm量产
  \item SK海力士CEO警告2027年将是史上最严重存储短缺
\end{itemize}
\begin{figure}[htbp]
\centering
\includegraphics[width=0.88\linewidth]{figures/fig_110.jpg}
\caption{Figure 1 - Figure 1: Liability-to-asset ratio of Chinese and global hyperscalers as of latest vs select Chinese real estate players and US telecom players in their respective cycle peaks Huawei Ascend 960 launch in 4Q — CXMT Shanghai fab e}
\end{figure}
\begin{figure}[htbp]
\centering
\includegraphics[width=0.88\linewidth]{figures/fig_111.jpg}
\caption{Figure 5 - Figure 5: China AI names' aggregate FTM P/S China AI ecosystem (≤US\textbackslash{}\$50bn) Figure 3: Expected timeline of select Chinese AI companies' production lift Figure 2: Chinese LLMs catching up with US peers}
\end{figure}
\begin{figure}[htbp]
\centering
\includegraphics[width=0.88\linewidth]{figures/fig_112.jpg}
\caption{Figure 5 - Figure 6: China AI names' aggregate standard deviation of FTM P/S}
\end{figure}
\begin{figure}[htbp]
\centering
\includegraphics[width=0.88\linewidth]{figures/fig_113.jpg}
\caption{Figure 4 - Figure 4: Expected timeline of select global AI companies' production lift Figure 6: China AI names' aggregate standard deviation of FTM P/S China AI ecosystem (≤US\textbackslash{}\$50bn) Figure 7: A-share velocity (5d rolling)}
\end{figure}
{\small\color{refgray}本节综合 1 篇研究；涉及 JPM。}
\newpage
\section{辅助信号｜战略咨询}
本板块整合波士顿咨询当日五份报告，聚焦AI代理重塑营销与治理、大型项目运营风险前置、以及绿色航运溢价缺口三大主题，揭示企业战略与执行之间的系统性裂缝。
\subsection{AI代理正在重塑品牌竞争与治理结构}
\textbf{AI代理作为消费决策的主动参与者，迫使品牌从叙事竞争转向可验证的表现竞争；同时，CEO与董事会在AI执行节奏上的认知错位，正成为企业转型的隐蔽阻力。}
\begin{itemize}
  \item 96\%的CMO相信AI正推动营销端到端转型，但仅31\%进入代理式执行阶段，多数营销领导者仍用旧方法应对新现实（R074）。
  \item AI代理实时评估品牌的服务、定价与投诉率，品牌承诺与运营现实之间的差距被放大；品牌管理从叙事游戏转向承诺与执行的一致性工程（R074）。
  \item 75\%的董事会成员认为自身AI知识不输同行，但超过1/3的CEO认为董事会高估了AI替代人类的能力，近40\%的CEO表示董事会缺乏对AI重塑增长战略的知情判断（R075）。
  \item 约60\%的CEO认为董事会正“催促AI转型”，这种由不确定性驱动的FOMO将治理层焦虑转化为对执行层的非理性压力，形成“速度-质量”困境（R075）。
  \item BCG建议董事会亲自试用AI工具、下沉一线考察、并参与跨行业交流，以建立对技术速度的共同感知，使董事会从监督者变为CEO在AI转型中的真正伙伴（R076）。
\end{itemize}
\begin{figure}[htbp]
\centering
\includegraphics[width=0.88\linewidth]{figures/fig_196.jpg}
\caption{EXHIBIT 1 - \#\# EXHIBIT 1 Although \$96\textbackslash{}\%\$ of CMOs Believe AI Is Driving End-to-End Marketing Transformation, Only \$31\textbackslash{}\%\$ Have Made Agentic Execution a Reality Pressure \$(\textbackslash{}\%)\$ 94\% of CMOs highlight increased expectations Question: How have your CEO's/executive team's expect}
\end{figure}
\begin{figure}[htbp]
\centering
\includegraphics[width=0.88\linewidth]{figures/fig_197.jpg}
\caption{Exhibit 2 - This creates two fundamental shifts. Because AI systems respond by combining products, services, and recommendations from multiple categories to address these expressed needs, they also reveal patterns of co-mingled demand that traditional research rarely unco}
\end{figure}
\subsection{大型项目运营风险在剪彩前数年已被锁定}
\textbf{全球资本项目规模超7万亿美元，但仅约十二分之一能按时按预算交付；运营失败并非开业日突发，而是战略定义与设计阶段早期假设累积的结果。}
\begin{itemize}
  \item 2026年全球交通、能源、工业和住宅领域资本研究预计超7万亿美元，94\%的大型项目业主计划维持或增加研究，但54\%的受访者报告交付挑战在过去一年加剧（R077）。
  \item 仅约十二分之一的大型项目能按时按预算交付，成本超支平均超过初始预算的60\%，工期延误常达一到两年（R077）。
  \item 运营不确定性在战略定义阶段即被嵌入：需求假设、运营模式与绩效目标一旦设定，后续设计施工中灵活性急剧收窄，纠正成本高昂（R077）。
  \item 项目进入运营后，累积风险迅速大规模显现：运营团队从执行转为应急，手动干预取代自动化，程序性变通替代系统可靠性（R077）。
  \item 运营可行性应从概念阶段作为明确要求进行不确定性测试，而非隐含结果；成功项目并非交付团队更优秀，而是运营思维嵌入了最早决策（R077）。
\end{itemize}
\begin{figure}[htbp]
\centering
\includegraphics[width=0.88\linewidth]{figures/fig_199.jpg}
\caption{Exhibit 1 - The consequences are severe. Congestion, system failures, safety incidents, reputational damage, regulatory scrutiny, and prolonged stabilization efforts follow. These breakdowns are often framed as execution failures at project delivery. Yet the fact that the}
\end{figure}
\begin{figure}[htbp]
\centering
\includegraphics[width=0.88\linewidth]{figures/fig_200.jpg}
\caption{EXHIBIT 1 - Once a project enters live operations, accumulated risk materializes quickly and at scale. Operational teams shift from execution to containment. Manual interventions replace automation. Procedural workarounds substitute for system reliability. Extraordinary e}
\end{figure}
\subsection{绿色航运溢价意愿上升，但预算缺口仍阻碍脱碳}
\textbf{货主愿意为绿色航运支付的溢价以每年超30\%的速度上升，但平均4\%的溢价远低于脱碳所需的10\%-15\%；承诺与资源之间的缺口是当前核心障碍。}
\begin{itemize}
  \item 全球航运业年排放约3\%的全球二氧化碳，若视为国家可排进前十；货主愿意支付的绿色溢价从2021年翻倍至2023年的4\%，但离脱碳所需的10\%-15\%仍有明显距离（R078）。
  \item 超过80\%的货主表示愿意为绿色航运支付溢价，但超过60\%的企业承诺减少Scope 3排放，其中不到一半配备了足够预算（R078）。
  \item 欧洲客户平均意愿支付溢价5.1\%，远高于南美洲的1\%；汽车和医疗保健行业溢价意愿最高，分别达6.1\%和6.0\%（R078）。
  \item 最有价值的付费群体是“欧洲食品饮料行业且有Scope 3承诺的客户”，其平均意愿支付溢价高达11\%，表明预算缺口集中在特定细分市场（R078）。
  \item 航运公司若只等客户主动出钱可能错过窗口期；真正有预算的客户群体非常集中，需针对性设计绿色服务与定价策略（R078）。
\end{itemize}
\begin{figure}[htbp]
\centering
\includegraphics[width=0.88\linewidth]{figures/fig_204.jpg}
\caption{Exhibit 1 - Exhibit 1 - Cargo Owners Are Willing to Pay Higher Premiums for Green Shipping, but Not Enough to Reach 2050 Decarbonization Goals WtP Premium How much are you willing to pay extra for carbon-neutral marine shipping, today? WtP Pr}
\end{figure}
\begin{figure}[htbp]
\centering
\includegraphics[width=0.88\linewidth]{figures/fig_205.jpg}
\caption{Exhibit 2 - Exhibit 2 - WtP Premiums Vary Widely by Industrial Sector and Level of Green Commitments Companies across all sectors willing to pay average premium of 4\% Avg WtP of companies by sector and commitment types (i.e., with or without Sc}
\end{figure}
\newpage
\section{辅助信号｜智库/国际机构}
本板块整合经合组织与世贸组织最新研究，聚焦三大主题：全球价值链升级路径与制度瓶颈、劳动力市场制度对生产率的影响、以及贸易救济与气候融资工具的结构性变化。
\subsection{全球价值链升级与制度瓶颈}
\textbf{经合组织多项研究揭示，发展中国家嵌入全球价值链后向高附加值环节移动的潜力巨大，但受制于制度碎片化、认证体系错位和政府中枢协调能力不足。}
\begin{itemize}
  \item 哥斯达黎加电子业上游度远低于OECD均值且自2015年持续下降，表明其生产活动集中于低附加值组装环节；若向价值链上游移动，出口、产出和实际工资均有望显著增长（R067）。
  \item 乌兹别克斯坦水安全投资环境评分显示，宏观框架虽有进步，但制度碎片化、成本回收能力弱、数据缺口严重，导致PPP项目落地困难，根源在现金流而非法律（R072）。
  \item OECD对FLA和WRAP两大劳工认证体系的评估指出，认证标准在“识别不确定性”环节覆盖薄弱，WRAP更是因侧重结果验证而非过程管理被判定“未对齐”，认证本身不能替代尽职调查（R065、R066）。
  \item 政府中枢能力建设报告基于六国实证，指出频繁机构重组带来管理成本，跨部门协调应优先选择委员会或任务导向模式，而非动辄合并部门（R068）。
\end{itemize}
\begin{figure}[htbp]
\centering
\includegraphics[width=0.88\linewidth]{figures/fig_170.jpg}
\caption{Figure 1 - \#\# 4. Empirical application: the case of Costa Rica Costa Rica's trade performance is associated with a strong capacity to attract large inflows of Foreign Direct Investment (FDI), which have been concentrated in sectors such as life sciences, light and advanc}
\end{figure}
\begin{figure}[htbp]
\centering
\includegraphics[width=0.88\linewidth]{figures/fig_191.jpg}
\caption{Figure 3 - \#\# 3.1. An overview of the OECD Scorecard for assessing the enabling environment for investment in water security The purpose of the OECD Scorecard is to pinpoint issues that hamper water security-related investments by public and private investors. The Scorec}
\end{figure}
\subsection{劳动力市场制度对生产率的隐性拖累}
\textbf{OECD首次跨国实证表明，竞业限制条款不仅压低工资、抑制跳槽，还通过阻碍劳动力向高效企业流动和知识扩散，系统性拖累整体生产率。}
\begin{itemize}
  \item 行业层面竞业限制覆盖率每上升10个百分点，整体劳动生产率水平下降约1.9\%，该效应在控制行业、国家、年份固定效应后依然稳健（R069）。
  \item 竞业限制条款越普遍的行业，生产率提升型劳动力再配置强度下降约17\%，即高效企业更难通过招聘获得人才，低效企业也不会因人才流失而被倒逼改善（R069）。
  \item 知识扩散同样受阻：竞业限制削弱了员工流动带来的技术和管理经验传播，落后企业追赶前沿的速度明显放缓（R069）。
  \item 该研究覆盖15个国家、近百万家企业级观测，首次将竞业限制与宏观生产率直接挂钩，指向保护商业利益的制度设计可能牺牲经济活力（R069）。
\end{itemize}
\begin{figure}[htbp]
\centering
\includegraphics[width=0.88\linewidth]{figures/fig_180.jpg}
\caption{Figure 1 - \#\# Box 1. Non-compete clauses and business entry Non-compete clauses directly prohibit employees from starting competing businesses. Though the practice is often justified by the need for incumbent firms to protect legitimate business interests and trade secre}
\end{figure}
\begin{figure}[htbp]
\centering
\includegraphics[width=0.88\linewidth]{figures/fig_181.jpg}
\caption{Figure 2 - \#\# Figure 2. Non-compete clauses are associated with weaker productivity-enhancing labour reallocation Non-compete clauses and the employment growth gap between high- and low-productivity firms Note: The employment growth gap is the difference in employment gr}
\end{figure}
\subsection{贸易救济与气候融资工具的结构性变化}
\textbf{美国对进口羊肉启动保障调查叠加碳信用市场从项目级向规模化转型，标志着贸易救济和气候融资工具正经历制度性升级，对全球供应链和资产叙事产生深远影响。}
\begin{itemize}
  \item 美国ITC于2026年7月13日应USTR请求启动羊肉保障调查，认定案件“异常复杂”，裁决截止日推迟至11月13日，最终报告可能跨越中期选举后的权力交接期（R073）。
  \item 调查聚焦羔羊肉，排除活羊和成年羊肉，核心证据链为进口量持续增长且价格显著低于国内产品，已导致企业退出、就业减少和产能利用率下降（R073）。
  \item 规模化碳信用模式正在成型：截至2025年超80个司法管辖区参与，但签发量仅占全球碳信用供给的6\%；其环境完整性取决于基线设定、额外性和碳泄漏追踪，若过关可能成为最具杠杆效应的气候融资工具（R071）。
  \item 捐赠国和企业已承诺约20亿美元用于规模化碳信用，主要投向司法管辖区REDD+项目，但全面实施案例仍有限，战略潜力大于当前规模（R071）。
\end{itemize}
\begin{figure}[htbp]
\centering
\includegraphics[width=0.88\linewidth]{figures/fig_186.jpg}
\caption{Figure 2 - \textbackslash{}- Article 4 outlines how governments set their mitigation strategies, targets and plans. Scaled-up crediting should support both the ambition and achievement of these efforts, and remain complementary to mitigation that would occur in the absence the creditin}
\end{figure}
\subsection{AI、技术与生产率}
\textbf{用于判断 AI 投资周期、生产率扩散和产业约束是否正在改变资产定价。}
\begin{itemize}
  \item 2022年底ChatGPT发布三年后，经合组织（OECD）最新报告揭示了一个看似矛盾的局面：英国本科生GenAI使用率从2024年初的66\%飙升至2026年的95\%，德国学生使用率从63\%跃升至90\%以上，法国从55\%升至82\%。在欧盟范围内，72\%的学生在过去三个月内使用过Ge...
\end{itemize}
\begin{figure}[htbp]
\centering
\includegraphics[width=0.88\linewidth]{figures/fig_185.jpg}
\caption{Figure 1 - \#\# Nearly all students and most academic staff have experience with GenAI In the three years since the public release of the first of a new generation of generative artificial intelligence (“GenAI”) tools in late 2022, they have developed from a novelty to bei}
\end{figure}
\section{结语}
后续需验证的数据与叙事节点包括：中国7月社零与固定资产投资数据能否确认消费动能实质改善；韩国央行8月会议加息概率及终端利率路径的进一步信号；台积电2026年资本开支执行节奏与AI芯片产能扩张进展；一线城市二手房价格走势（NBS与中原指数分歧的收敛方向）；印度即时零售行业补贴削减节奏与单店订单密度变化。
\newpage
\section{覆盖概览}
投行/券商 64 篇；战略咨询 5 篇；智库/国际机构 9 篇。\par
投行覆盖：GS 19篇 / JPM 15篇 / MS 15篇 / Bernstein 5篇 / Citi 4篇 / NOM 3篇 / BARC 1篇 / DB 1篇 / HSBC 1篇\par
\section{Disclaimer}
{\scriptsize\color{refgray}
This community is a paid learning and information-sharing community created for general educational purposes only. The membership fee is charged solely for access to community discussions, educational content, organization, and operational costs. It is not an advisory fee, management fee, performance fee, brokerage fee, or compensation for personalized financial, investment, legal, tax, or professional advice.\par
I participate and share information solely as an individual and for educational discussion among community members. I am not acting as a financial adviser, investment adviser, broker, dealer, portfolio manager, legal adviser, tax adviser, accountant, fiduciary, or any other licensed professional adviser. Nothing shared in this community constitutes, and should not be interpreted as, financial advice, investment advice, legal advice, tax advice, a recommendation, solicitation, offer, endorsement, instruction, or guarantee to buy, sell, hold, short, trade, or invest in any security, cryptocurrency, fund, derivative, strategy, or other financial product.\par
All content shared in this community, including messages, discussions, links, files, charts, examples, market commentary, watchlists, AI-generated or AI-polished summaries, and third-party materials, is general, impersonal, and educational in nature. It does not take into account any individual's financial situation, investment objectives, risk tolerance, investment horizon, liquidity needs, tax status, legal restrictions, or suitability requirements. No content should be relied upon as suitable for any specific person, account, portfolio, or financial decision.\par
Information shared in this community may be incomplete, inaccurate, outdated, speculative, AI-generated, AI-edited, or based on third-party internet sources that have not been independently verified. I make no representation or warranty as to the accuracy, completeness, timeliness, reliability, or suitability of any information shared.\par
Financial markets involve substantial risk, including the possible loss of principal. Past performance, historical data, hypothetical examples, simulated results, backtests, personal opinions, or educational examples do not guarantee future results. Each member is solely responsible for conducting their own research, exercising independent judgment, and making their own decisions. Before making any financial, investment, legal, or tax decision, members should consult qualified and licensed professionals.\par
Participation in this community, payment of any membership fee, or communication with me or other members does not create any adviser-client, fiduciary, professional, contractual, agency, partnership, employment, or other advisory relationship. By joining or participating in this community, you acknowledge and agree that you are solely responsible for your own decisions, actions, transactions, profits, losses, risks, and consequences, and that I shall not be liable for any direct, indirect, incidental, consequential, financial, legal, tax, or other loss or damage arising from or related to any content, discussion, information, or material shared in this community.\par
}
\newpage
\begin{center}
{\Large 更多详情报告kcdesk.com}\par\vspace{0.8cm}
\includegraphics[width=0.58\linewidth]{\detokenize{../../prompts/zsxq_img.jpg}}
\end{center}
\end{document}
